\section[Hyperbolic deformation and the time component inequality]{Removal of the second fundamental form on an asymptotically hyperbolic end}
\label{sec:ah-deformation}

All norms and differentiated estimates in this section are measured in the
hyperbolic reference metric on balls of fixed hyperbolic radius, unless
another convention is stated.  Thus \(O_j(r^{-a})\) includes covariant
derivatives through order \(j\), without a loss of a power of \(r\).
The letter \(r\) is extended to a positive smooth function on the compact
part of the manifold.  An enclosing cut always means the entire intrinsic
boundary of a connected exterior containing the designated end.

\begin{theorem}[The time component inequality]\label{thm:ah-component}
Let \((\Omega,g,K)\) satisfy the regularity, completeness, one-end,
dominant energy, weighted integrability, and future boundary hypotheses
of the initial-data class, with
\[
 g-b\in C^{2,\alpha}_{\tau_{\rm data}},\qquad
 K\in C^{1,\alpha}_{\tau_{\rm data}},\qquad
 \frac32<\tau_{\rm data}<3 .
\]
Suppose the four metric fluxes in the designated chart are finite.
There is no causal or sign assumption on their vector in this theorem.
Then
\begin{equation}
 p_0(g)\ge
 \sqrt{\frac{A_{\min}(S;g)}{16\pi}}
 \left(1+\frac{A_{\min}(S;g)}{4\pi}\right).
 \label{eq:ah-component-bound}
\end{equation}
The tensor \(K\) may have arbitrary trace, and the boundary may have
arbitrarily many components.
\end{theorem}

We shall prove this by constructing comparison metrics with scalar
curvature at least \(-6\).  The time-symmetric case of
Theorem~\ref{thm:str-one-sign}, whose proof uses
Theorem~\ref{thm:fl-deformation}, will then apply in this same chart.
The continuation and estimates below use the geometric and elliptic
lemmas proved in the preceding sections.

\subsection{Strict approximation with improved tensor decay}
\label{subsec:ah-preparation}

Write
\[
 M= \operatorname{div}_g\bigl(K-(\tr_gK)g\bigr)=8\pi J,
 \qquad
 Q_g(K)=(\tr_gK)^2-|K|_g^2,
 \qquad
 M_d=8\pi(\mu-|J|_g).
\]
Choose once and for all
\begin{equation}\label{eq:ah-preparation-exponents}
 2<\kappa<1+\sqrt3,\qquad
 \frac32<\tau'<\min\{\kappa,\tau_{\rm data}\},\qquad
 0<\delta<\min\{1,2\kappa-4,2\tau'-3\}.
\end{equation}
These choices are possible even when \(\tau_{\rm data}\) is arbitrarily
close to \(3/2\).

\begin{lemma}[Strict approximation]\label{lem:ah-preparation}
There are smooth data \((g_j,K_j)\) converging to \((g,K)\) on compact
sets, whose metric bilinear ratios converge uniformly to one, such that
all four metric mass components converge and
\begin{equation}\label{eq:ah-prepared-data}
 \begin{split}
 &8\pi(\mu_j-|J_j|_{g_j})\ge c_j(1+r)^{-3-\delta},
 \qquad H_{g_j}+\tr_TK_j<0,\\
 &g_j-b=O_\infty(r^{-\kappa}),\qquad
 K_j=O_\infty(r^{-\kappa}),\qquad
 R_{g_j}+6=O_\infty(r^{-3-\delta})
 \end{split}
\end{equation}
for positive constants \(c_j\).  In particular
\(A_{\min}(S;g_j)\to A_{\min}(S;g)\).
\end{lemma}

\begin{proof}
Let \(\chi_R\) equal one on \(r\le R\) and zero on \(r\ge2R\), with
uniform differentiated bounds in \(\log r\).  Interpolate the metric to
\(b\) by \(q_R=b+\chi_R(g-b)\), keeping the original metric on the
compact part.  Let \(P=K-(\tr_gK)g\).  On the exterior of a fixed
sufficiently large sphere \(r=L\), solve
\begin{equation}\label{eq:ah-vector-corrector}
 (\Delta_{q_R}+\Ric_{q_R})X_R
 =\chi_R\operatorname{div}_gP-
       \operatorname{div}_{q_R}(\chi_RP),\qquad
 X_R|_{r=L}=0,\qquad X_R\longrightarrow0 .
\end{equation}
The vector Laplacian acts on one-forms.  On this fixed exterior the
Ricci endomorphism is bounded above by a negative multiple of the
identity, uniformly for large \(R\).  The source has
\(C^{0,\alpha}_{\tau'}\) norm \(O(R^{\tau'-\tau_{\rm data}})\).
Kato's Inequality compares \(|X_R|\) with the scalar operator
\(-\Delta+2+o(1)\).  For a radial power,
\[
 (-\Delta_b+2)r^{-a}
 =\bigl(2+2a-a^2+O(r^{-2})\bigr)r^{-a}.
\]
The coefficient is positive for \(a<1+\sqrt3\), so \(r^{-\tau'}\)
is a common supersolution after enlarging \(L\).

For clarity, existence in this argument follows on finite annuli from
the coercivity of \(-\Delta-\Ric\) with zero Dirichlet values.
The maximum comparison just given and local boundary estimates are
independent of the outer annulus radius.  Dirichlet exhaustion,
followed by local Schauder estimates, gives
\[
 X_R=O_2(\varepsilon_Rr^{-\tau'}),\qquad
 \varepsilon_R=R^{\tau'-\tau_{\rm data}}\longrightarrow0.
\]
On \(r>2R\) the metric is exactly \(b\) and the source vanishes.
The same comparison with \(r^{-\kappa}\), and differentiation on
hyperbolic balls, gives \(X_R=O_\infty(r^{-\kappa})\) there, with
constants allowed to depend on \(R\).

Choose a fixed cutoff equal to zero near \(r=L\) and one beyond a
slightly larger sphere, and add the cutoff of
\[
 2\operatorname{sym}\nabla X_R-(\operatorname{div}X_R)q_R
\]
to \(\chi_RP\).  Denote the resulting stress by \(P_R\), and let \(K_R\)
be its inverse trace reversal in \(q_R\).  The identity
\[
 \operatorname{div}_{q_R}
 \bigl(2\operatorname{sym}\nabla X_R
             -(\operatorname{div}X_R)q_R\bigr)
 =(\Delta_{q_R}+\Ric_{q_R})X_R
\]
shows that
\[
 \operatorname{div}_{q_R}P_R
 =\chi_R\operatorname{div}_gP+E_R ,
\]
where \(E_R\) is supported in a fixed annulus and tends to zero in
every needed compact norm.  The data are unchanged near \(S\).
The family \(K_R\) is bounded by a common smooth weight of order
\(-\tau'\), and its exact tail has order \(-\kappa\).

Here is the error estimate needed for the energy condition.  Let
\(w_0>0\) be smooth and equal to \(r^{-3-\delta}\) on the tail.
There is a number \(e_R\to0\) such that
\begin{equation}\label{eq:ah-cutoff-dec}
 \chi_R(R_g+6)+Q_{q_R}(K_R)
       -2|\operatorname{div}_{q_R}P_R|_{q_R}
 \ge -e_Rw_0 .
\end{equation}
Indeed the expression with \(Q_{q_R}(K_R)\) replaced by
\(\chi_RQ_g(K)\), and with the last norm in \(g\), is nonnegative
by the original dominant energy condition.  Its errors are the
fixed-annulus term \(E_R\), products involving \(X_R\) and \(K\),
the quadratic cutoff error \((\chi_R^2-\chi_R)Q_g(K)\), and the
change of norm between \(g\) and \(q_R\).  The noncompact terms
are bounded by a coefficient tending to zero times
\(r^{-2\tau'}\).  Since \(2\tau'>3+\delta\), these are
\(o(1)w_0\).  The fixed-annulus terms have the same conclusion by
compact positivity of \(w_0\).  This argument uses no extra
pointwise decay of the original density \(\mu\).

Let \(D\ge |K_R|_{q_R}\) be a common smooth positive weight of order
\(-\tau'\), and let \(w_2>0\) be a smooth weight equal to \(r^{-\tau'}\)
far out.  Choose \(\eta_R\to0\) sufficiently slowly that
\(e_R+\varepsilon_R=o(\eta_R)\).  Solve
\begin{equation}\label{eq:ah-conformal-corrector}
 \begin{split}
 (-\Delta_{q_R}+3)v_R={}&
 \frac18\{\chi_R(R_g+6)-(R_{q_R}+6)\}\\
 &+D\sqrt{|dv_R|^2+\eta_R^2w_2^2}+\eta_Rw_0,\\
 \partial_\nu v_R|_S={}&-\eta_R,\qquad v_R\longrightarrow0 .
 \end{split}
\end{equation}
The scalar forcing satisfies
\[
 \left\|\tfrac18\{\chi_R(R_g+6)-(R_{q_R}+6)\}\right\|_{C^{0,\alpha}_{\tau'}}
 \le C\varepsilon_R=o(\eta_R),
\]
by the differentiated cutoff estimates above.
Here and below \(\nu\) at an inner boundary points into the exterior.
This is a scalar equation with a bounded first-order drift and a
positive zeroth-order coefficient.  A complete continuation argument
is as follows.  On a finite truncation impose zero outer values and
multiply the right side and the inner normal derivative by a
parameter \(a\in[0,1]\). Subtract \(a\eta_R\) times a fixed collar
function with normal derivative \(-1\); the remaining function has
homogeneous Neumann data.  At its extrema the gradient term is
bounded by the fixed collar derivatives and the regularizing
source.  The Neumann maximum principle therefore gives a uniform
\(O(\eta_R)\) bound.  On the end, \(r^{-\tau'}\) is a supersolution
for the equation with any of its bounded decaying drifts, because
\(3+2\tau'-(\tau')^2>0\).  Hence
\[
 v_R=O_2(\eta_Rr^{-\tau'}).
\]
The estimates of Lemma~\ref{lem:linear-separated-faces}, first
absorbing the bounded first-order term by interpolation and then
applying the one-sided Schauder estimate, give closedness
of continuation.  Its linearization has positive zeroth order and
the same mixed boundary conditions, so the maximum principle and
linear elliptic theory give openness.  Exhaustion now gives the
displayed solution.  On the exact tail the sources decay faster
than \(r^{-\kappa}\); comparison improves the bound to that power.
The regularizer \(w_2\) bounds all differentiated coefficients on
rescaled balls.  Bootstrap gives the stated estimates of all
orders at fixed \(R\).

Set \(g_R'=e^{4v_R}q_R\), \(K_R'=e^{2v_R}K_R\).  Direct conformal
calculation gives
\begin{align}
 8\pi J_R'&=e^{-2v_R}
  \{\,\operatorname{div}_{q_R}P_R+4K_R(\nabla v_R,\cdot)\,\},
       \label{eq:ah-conformal-momentum}\\
 e^{4v_R}(R_{g_R'}+6)&=\chi_R(R_g+6)
  +8D\sqrt{|dv_R|^2+\eta_R^2w_2^2}+8\eta_Rw_0
  +6(e^{4v_R}-1-4v_R)-8|dv_R|^2 .
       \label{eq:ah-conformal-density}
\end{align}
The norm in the first equation contributes another factor \(e^{-2v_R}\).
Equations~\eqref{eq:ah-cutoff-dec}--\eqref{eq:ah-conformal-density},
\(D\ge|K_R|\), and \(2\tau'>3+\delta\) imply strict dominant energy
with a positive multiple of \(\eta_Rw_0\).  At the boundary
\[
 \theta_+(g_R',K_R')
 =e^{-2v_R}\{\theta_+(g,K)+4\partial_\nu v_R\}
 =e^{-2v_R}\{\theta_+(g,K)-4\eta_R\}<0 .
\]
Moreover, the four terms following \(\chi_R(R_g+6)\) in
Equation~\eqref{eq:ah-conformal-density} are \(O(\eta_R)w_0\).

We include the mass argument because uniform metric convergence alone
would not prove it.  The original constraint and the weighted
integrability assumption imply
\(\int V_0|R_g+6|\,dV_g<\infty\).
For any static potential \(V_a\), its scalar-linearization flux
has divergence \(V_aDR_b(g-b)\).  The difference between
\(DR_b(g-b)\) and \(R_g+6\) is quadratic in the metric error and its
first two derivatives.  Its weighted radial integral is bounded by
\[
 C\int^\infty r^{2-2\tau'}\,dr<\infty .
\]
Equation~\eqref{eq:ah-conformal-density} makes the curvature tail
uniformly small, and the quadratic error has the same uniform
integrable bound.  Compare the fluxes on a fixed large sphere,
where the data converge, and then send that sphere to infinity.
Since \(|V_a|+|dV_a|\le C V_0\), this proves convergence of all
four components.  It also proves their existence for each
approximating metric.  The improved tail bounds and the constraints
give all the weighted integrability requirements for the new data.
If the original data have only the stated finite regularity, the
constructed data are already smooth outside a compact set: there
the vector equation is homogeneous on \(b\), and all coefficients
and sources of the scalar equation are smooth after bootstrap.
On the remaining compact set, smooth the metric in \(C^2\) and the
tensor in \(C^1\), with a cutoff leaving that smooth tail unchanged.
The strict energy margin and strict boundary expansion have positive
minima in absolute value on their respective compact sets. Continuity
of scalar curvature in \(C^2\), momentum divergence in \(C^1\), and
boundary expansion in these norms therefore preserves both
inequalities if the smoothing error is sufficiently small.
Choose that error to tend to zero with \(R\). This gives smooth
approximants with unchanged asymptotic estimates and mass components.

Finally, uniform bilinear ratio convergence bounds every cut area
between factors tending to one, and therefore bounds their infima
by the same factors.
\end{proof}

It suffices henceforth to prove Theorem~\ref{thm:ah-component} for
fixed data satisfying Equation~\eqref{eq:ah-prepared-data}.
All constants from this preparation are fixed before any parameter
of the following deformation is chosen.

\subsection{The compact domain, the lapse, and the profiles}
\label{subsec:ah-profiles}

Coordinate spheres have both expansions tending to \(2\).  The
maximum-radius comparison therefore confines all bounding regions
at negative thresholds near zero to a common compact set.  Apply
Lemmas~\ref{lem:fl-mots-barriers},
\ref{lem:fl-right-continuity}, and
\ref{lem:fl-frontier-collars} as follows.  First take the full black
region at a fixed right-continuity threshold \(c_b<0\), using the
given boundary as required inner boundary.  In its complement take
the full white region for \(-K\), at a right-continuity threshold
\(c_w<0\); the reversed black faces are strict outer barriers, and
this second region may be empty.  Retain the closed component
toward infinity after deleting both regions.  Its boundary consists
of entire black and white faces.  Every full cut in this smaller
exterior is a full cut in the original one.

The compact geometric lemmas apply without a cosmological or energy
assumption: their hypotheses are smooth compact geometry, the stated
strict outer barriers, and the tensor threshold shift
\(K\mapsto K-cg/2\).  Their generalized-support conclusion is
Lemma~\ref{lem:fl-support-enclosure}; it will be used below only on
fixed compact sets with those same barrier conventions.

Choose a smooth compactly supported offset \(C\), with
\(C_w\le C\le C_b\), such that on the outward leaf collars
\[
 C=C_b-k_Bs\quad\hbox{on black collars},\qquad
 C=C_w+k_Bs\quad\hbox{on white collars},
\]
where \(C_b>0>C_w\) and \(k_B>0\).  Set
\begin{equation}\label{eq:ah-face-heights}
 b_-=c_b-C_b<0,\qquad b_+=-c_w-C_w>0 .
\end{equation}
The thresholds and these offsets will be made small in the order
specified in Lemma~\ref{lem:ah-heights}.

\begin{lemma}[A spatial lapse]\label{lem:ah-lapse}
There is a positive smooth function \(\lambda\), constant on a fixed
large compact region containing all boundary collars, and equal to
\(\sqrt{A_0^2+r^2}\) far out, such that, with \(s=d\log\lambda\),
\begin{equation}\label{eq:ah-lapse-inequalities}
 |s|\le1,\qquad
 E:=\frac{\Delta\lambda-\Hess\lambda(w,w)}{\lambda}
           +v|s|_{A_d}^2\le3 .
\end{equation}
Here \(w\) is any vector of length less than one, \(v=|w|^2\), and
\(A_d=I-w\otimes w\).  On the end, also
\[
 E\le E_0:=\Delta\lambda/\lambda\le3,\qquad
 E_0=3+O(r^{-2}).
\]
A fixed smooth positive \(\rho\), radial of order \(-3-\delta\) on
the end, may be chosen with \(10\rho<M_d\), and the lapse can be
chosen so that \(|s||K|\le c_*\sqrt\rho\) for any prescribed
sufficiently small \(c_*>0\).
\end{lemma}

\begin{proof}
First smoothly replace \(r^2\) by a constant on a sufficiently large
compact set.  On a fixed compact region, taking \(A_0\) large makes
all derivatives of the logarithm of the resulting square root
small, proving Equation~\eqref{eq:ah-lapse-inequalities} there.
For the exact hyperbolic metric on the tail, the tangential and
radial Hessian eigenvalues divided by \(\lambda\) are
\[
 \frac{1+r^2}{A_0^2+r^2},
 \qquad
 1-\frac{A_0^4-A_0^2}{(A_0^2+r^2)^2}.
\]
Each is at least
\[
 |d\log\lambda|_b^2
 =\frac{r^2(1+r^2)}{(A_0^2+r^2)^2}.
\]
Their gaps are positive multiples of \(r^{-2}\) for \(A_0>1\).
Also
\[
 \frac{\Delta_b\lambda}{\lambda}-3
 =-\frac{(A_0^2-1)(3A_0^2+2r^2)}
           {(A_0^2+r^2)^2}.
\]
The errors from \(g-b=O_2(r^{-\kappa})\) are smaller than these gaps
because \(\kappa>2\).  This proves the assertions after enlarging
the fixed compact transition.  The last smallness condition follows
by making \(s=0\) throughout a sufficiently large compact set:
on its complement \(r^{-\kappa}=o(r^{-(3+\delta)/2})\).
\end{proof}

Choose \(\beta,q\) with
\begin{equation}\label{eq:ah-beta-q}
 2+\delta/2<\beta<\min\{3,\kappa\},
 \qquad q>\max\{3+\delta,2\beta\}.
\end{equation}
The constants \(\delta_0>0\), \(P_0\), and \(B_0\) below are chosen
successively in the floor and absorption arguments.  For a large
integer \(N\), put
\begin{equation}\label{eq:ah-profiles}
 \epsilon=\frac{B_0\log N}{N},\qquad
 \ell=N^{-B_0},\qquad \tau=\ell^{5/4},\qquad R=N^a .
\end{equation}
Choose the fixed exponent \(a\) so large that
\begin{equation}\label{eq:ah-R-conditions}
 a(2\beta-4-\delta)>\frac54B_0,\qquad
 a(q-1)+1>\frac54B_0.
\end{equation}
Let \(A_R=1\) before \(R\), let \(A_R\asymp1+r/R\), and let it be
proportional to \(r\) beyond \(2R\).  Its logarithmic derivative in
\(\log r\) is nondecreasing from zero to one.  Let \(n=N\) on a
fixed compact region, let \(n\asymp N(1+r)^q\) up to \(R\), and
cap it smoothly to \(n_\infty\asymp NR^q\) past \(2R\).  The
positive-order derivatives of \(\log n\) are uniformly bounded.
The implicit comparison constants may depend on the fixed inner
radius where \(n\) begins to increase.  Equations
\eqref{eq:ah-R-conditions} imply
\begin{equation}\label{eq:ah-small-ratios}
 \frac{R^{4+\delta-2\beta}}{\tau}\longrightarrow0,\qquad
 \frac{R}{n_\infty\tau}\longrightarrow0 .
\end{equation}
All these profiles are held fixed while an independent outer
truncation radius tends to infinity.

Let \(p_*\) be smooth and nonincreasing, equal to one on
\((-\infty,0]\) and zero on \([1,\infty)\).  Define
\[
 \begin{gathered}
 l_0(t,x)=\exp\!\left(\int_0^{n(x)t}p_*(z)\,dz\right),\qquad
 l=\lambda l_0,\\
 p=\partial_t\log l_0=np_*(nt),\qquad
 A=d_x\log l=s+pt\,d\log n .
 \end{gathered}
\]
Spatial derivatives carrying the subscript \(x\) keep \(t\) fixed.
For \(t\le0\), \(l_0=e^{nt}\); for \(t\ge1/n\), \(l_0\) is a
constant independent of \(n\).

\subsection{Variables and the physical equations}
\label{subsec:ah-equations}

For functions \(f,t\), define
\[
 \begin{gathered}
 \sigma=|df|,\quad D=1+l^2\sigma^2,\quad d=D^{-1},
 \quad w=\frac{l\nabla f}{\sqrt D}=ae,\quad
 a=|w|,\quad v=a^2,\quad e=\nabla f/\sigma,\\
 A_j=I+(j-1)e\otimes e,\qquad
 \chi=\frac{4d}{4+pv},\qquad
 H^f=\frac{l\Hess f}{\sqrt D},\\
 S=K+H^f+w\otimes A+A\otimes w,\qquad
 h=\tau A_Rf,\qquad
 L=e^{2t}l_0,\quad u=\lambda L\sqrt D,\quad
 Z=t+\frac18\log D .
 \end{gathered}
\]
All actual coefficients extend smoothly across \(\sigma=0\);
for example \(A_\chi=I-(4+p)w\otimes w/(4+pv)\).
For fixed \(x,df\), the derivative \(\partial Z/\partial t\) is
\((4+pv)/4>0\), and the map is onto \(\mathbb R\).
Thus \(t=t(x,Z,df)\) is defined without imposing a floor.

Decompose \(S\) into its transverse block \(P\), mixed block \(M\),
and axial entry \(q_1\).  Write \(x_t=\partial_et\) and
\(y=\nabla_\perp t\).  In expressions involving the trace equation,
\(F=\tr_{A_\chi}S\), so \(\tr P=F-\chi q_1\).  Put
\begin{equation}\label{eq:ah-T}
 \begin{split}
 \mathcal T={}&\frac12|P^{\rm tf}|^2
 +|M+(p+1)ay|^2+[4(p+1)-v]|y|^2\\
 &+4(p+1)d\left(x_t+\frac{\chi a q_1}{4d}\right)^2
 +\frac34\left(F-\frac{\chi q_1}{3}\right)^2
 +\frac{4d(9-d)}{3(4+pv)^2}q_1^2 .
 \end{split}
\end{equation}
Let \(d_0(nt)\) equal \(1/2\) for \(nt\ge0\), equal a fixed
sufficiently small \(\delta_0>0\) for \(nt\le-1\), and lie between
those values.  Let \(\rho_0>0\) have order \(-3-\delta\), and let
\(m_0(t,x)=m(n(x)(t+\epsilon))\), where \(m=1\) on
\((-\infty,0]\), \(m=0\) on \([1,\infty)\), and \(0\le m\le1\).
The physical equations are
\begin{equation}\label{eq:ah-system}
 \begin{split}
 \tr_{A_\chi}S&=F=h+C,\\
 \operatorname{div}V&=u\Xi,\qquad
 V=uA_\chi\{4\nabla Z+K(w,\cdot)+A(w)w\},\\
 \Xi&=3(e^{4t}-1)+4\langle A,dt\rangle_{A_d}
 +d_0(\mathcal T+\tau A_Ra\sigma)+\rho
 -m_0n^{P_0}(\rho_0+v).
 \end{split}
\end{equation}
On black faces impose \(h=b_-\), on white faces \(h=b_+\), and
on every inner face impose
\begin{equation}\label{eq:ah-inner-flux}
 \frac{V_\nu}{u}
 =\mathcal B:=-H+w_\nu(F-P_B)-N_Bd,\qquad
 P_B=\tr_TK,\qquad N_B>1+\sup|H|.
\end{equation}
On an outer truncation sphere impose \(f=Z=0\).

The comparison metrics are
\[
 \bar g=g+l^2df^2,\qquad \hat g=e^{4t}\bar g.
\]
The exact identity underlying the equations is
\begin{equation}\label{eq:ah-curvature}
\begin{aligned}
 \frac12e^{4t}(R_{\hat g}+6)
 &=8\pi(\mu+J(w))+3(e^{4t}-1)
  +4\langle A,dt\rangle_{A_d}
  +\mathcal T\\
 &\quad+w(F)-F\tr K-u^{-1}\operatorname{div}V .
\end{aligned}
\end{equation}
Its full derivation, together with the weighted primitive and the
metric flux normalization, is given next.

\subsection{Exact identities and the metric mass flux}
\label{sec:ah-algebra}

We give the differential identities used in the deformation, including
their dependence on the spatially varying lapse profile.  The argument
in this subsection is local until the asymptotic assumptions in
Proposition~\ref{prop:ah-alg-mass-passage} are imposed.  In particular,
none of the identities requires a solution of the deformation equations
to have been constructed in advance.

\paragraph{Notation.}
All gradients, contractions, and divergences in the differential
identities below refer to \(g\).  We identify vectors and covectors using
\(g\) when they occur in the same formula.  Let \(t,f\) be smooth,
let \(l(t,x)>0\), and write
\[
 p=\partial_t\log l,\qquad
 A=\mathrm d_x\log l,\qquad
 B=\mathrm d\log l=A+p\,\mathrm dt.
\]
Here \(\mathrm d_x\) holds \(t\) fixed, whereas \(\mathrm d\) differentiates
the composition \(x\mapsto l(t(x),x)\). Use \(l=\lambda l_0\)
from Section~\ref{subsec:ah-profiles} and \(L=e^{2t}l_0\)
from Section~\ref{subsec:ah-equations}. Set
\begin{gather*}
 D=1+l^2|\mathrm df|^2,\qquad d=D^{-1},\qquad
 w=\frac{l\nabla f}{\sqrt D},\qquad v=|w|^2=1-d,\qquad a=\sqrt v,\\
 H^f=\frac{l\nabla^2 f}{\sqrt D},\qquad
 \overline g=g+l^2\mathrm df\otimes\mathrm df,\qquad
 \widehat g=e^{4t}\overline g,\\
 U=l\sqrt D,\qquad u=e^{2t}U=\lambda L\sqrt D,\qquad
 Z=t+\frac18\log D .
\end{gather*}
At a point where \(a>0\), put \(e=w/a\), and define
\[
 A_j=I+(j-1)e\otimes e,\qquad
 \chi=\frac{4d}{4+pv},\qquad
 S=K+H^f+w\otimes A+A\otimes w,\qquad
 F=\operatorname{tr}_{A_\chi}S.
\]
Thus \(A_d=I-w\otimes w\), which is the inverse of \(\overline g\)
relative to \(g\).  We assume \(p\ge0\), as holds for the profiles used
here.  The formulas involving \(A_\chi\) extend smoothly across
\(a=0\), since
\[
 A_\chi=I-\frac{p+4}{4+pv}\,w\otimes w.
\]
For a symmetric tensor \(T\), let
\[
 p_T=T-(\operatorname{tr}_gT)g,\qquad
 [T_1,T_2]=\langle T_1,T_2\rangle
          -(\operatorname{tr}_gT_1)(\operatorname{tr}_gT_2).
\]
We use the constraint normalization
\[
 16\pi\mu=R_g+6+(\operatorname{tr}_gK)^2-|K|^2,\qquad
 8\pi J=\operatorname{div}_g p_K.
\]

Decompose \(S\) into its \(e^\perp\)-block \(P\), mixed block \(M\),
and axial entry \(q_1=S(e,e)\), and put
\(x_t=\partial_e t\), \(y=(\mathrm dt)_{e^\perp}\).
The quadratic form of Equation~\eqref{eq:ah-T} is
\begin{align}
 \mathcal T={}&\frac12|P^{\mathrm{tf}}|^2
       +|M+(p+1)a y|^2+[4(p+1)-v]|y|^2 \notag\\
 &+4(p+1)d\left(x_t+\frac{\chi a q_1}{4d}\right)^2
       +\frac34\left(F-\frac{\chi q_1}{3}\right)^2
       +\frac{4d(9-d)}{3(4+pv)^2}q_1^2 .
 \label{eq:ah-alg-T}
\end{align}
Although this expression uses an axis where \(a>0\), its sum is
direction independent at \(a=0\), as will also follow from the
coordinate-free expression in the proof below.

\begin{proposition}[Curvature identity]
\label{prop:ah-alg-curvature}
With the preceding definitions, put
\begin{equation}
 V=uA_\chi\bigl(4\nabla Z+K(w,\cdot)^\sharp+A(w)w\bigr).
 \label{eq:ah-alg-V}
\end{equation}
Then the following identity holds pointwise:
\begin{align}
 \frac12e^{4t}(R_{\widehat g}+6)
 ={}&8\pi\bigl(\mu+J(w)\bigr)+3(e^{4t}-1)
       +4\langle A,\mathrm dt\rangle_{A_d}
       +\mathcal T \notag\\
 &+w(F)-F\operatorname{tr}_gK-u^{-1}\operatorname{div}_g V .
 \label{eq:ah-alg-curvature-identity}
\end{align}
Here \(F\) denotes the actual trace
\(\operatorname{tr}_{A_\chi}S\); therefore the identity applies to
any prescribed trace equation by substituting that prescription for
\(F\).
\end{proposition}

\begin{proof}
We first derive the graph identity, keeping the full spatial
dependence of \(l\).  On a product with coordinate \(s\), consider the
stationary Lorentz metric
\[
 \mathbf g=-l^2(\mathrm ds-\mathrm df)^2+\overline g
           =-l^2\mathrm ds^2+2l^2\mathrm ds\,\mathrm df+g .
\]
The \(s\)-slices have induced metric \(g\), shift
\(\beta=l^2\nabla f=Uw\), and lapse \(U=l\sqrt D\).  For the future
normal \(U^{-1}(\partial_s-\beta)\), with the convention
\(k(X,Y)=\mathbf g(\mathbf\nabla_X n,Y)\), stationarity gives
\[
 k=-\frac{1}{2U}\mathcal L_\beta g
   =-H^f-w\otimes B-B\otimes w .
\]
The contracted Gauss equation and the normal Ricci equation for these
stationary slices give
\begin{equation}
 U R_{\mathbf g}
 =U\bigl(R_g+[k,k]\bigr)
   -2\operatorname{div}_g\bigl(\nabla U+(\operatorname{tr}_g k)\beta\bigr).
 \label{eq:ah-alg-stationary}
\end{equation}
For completeness, this sign can be fixed directly by combining
\[
 R_{\mathbf g}=R_g+(\operatorname{tr}k)^2-|k|^2
                   -2\operatorname{Ric}_{\mathbf g}(n,n)
\]
with the normal expansion identity
\[
 \operatorname{Ric}_{\mathbf g}(n,n)
 =-n(\operatorname{tr}k)-|k|^2
       +U^{-1}\Delta_g U .
\]
Indeed \(n(\operatorname{tr}k)=-U^{-1}\beta(\operatorname{tr}k)\)
by stationarity, and
\(\operatorname{div}_g\beta=-U\operatorname{tr}_g k\).
Substitution is exactly
Equation~\eqref{eq:ah-alg-stationary}.

In the coordinate \(s-f\), the same spacetime metric is
\(-l^2\mathrm d(s-f)^2+\overline g\), so
\[
 R_{\mathbf g}=R_{\overline g}-2l^{-1}\overline\Delta l .
\]
The inverse metric and volume element of \(\overline g\) give
\[
 \overline\Delta l
   =D^{-1/2}\operatorname{div}_g(\sqrt D\,A_d\nabla l).
\]
Differentiating \(U\) also gives
\[
 \nabla U-\frac Ul A_d\nabla l=-k(\beta,\cdot)^\sharp.
\]
It follows from Equation~\eqref{eq:ah-alg-stationary} that
\[
 U R_{\overline g}
  =U\bigl(R_g+[k,k]\bigr)+2\operatorname{div}_g(p_k\beta).
\]
Let
\[
 Q=K-k=S+p(w\otimes\mathrm dt+\mathrm dt\otimes w).
\]
Since the symmetric part of \(\nabla\beta\) is \(-Uk\),
\begin{align*}
 \frac2U\operatorname{div}_g(p_K\beta)
  &=2(\operatorname{div}_g p_K)(w)
       +\frac2U\langle p_K,\nabla\beta\rangle\\
  &=16\pi J(w)-2[K,k].
\end{align*}
Using \(p_k=p_K-p_Q\) and
\(R_g=16\pi\mu-6+[K,K]\), we obtain
\begin{equation}
 R_{\overline g}
  =16\pi\bigl(\mu+J(w)\bigr)-6+[Q,Q]
       -2U^{-1}\operatorname{div}_g(Up_Qw).
 \label{eq:ah-alg-graph-curvature}
\end{equation}

The conformal scalar-curvature formula in three dimensions is
\[
 e^{4t}R_{\widehat g}
   =R_{\overline g}-8\overline\Delta t
                         -8|\mathrm dt|_{\overline g}^2 .
\]
Replacing \(U\) by \(u=e^{2t}U\) in the first divergence in
Equation~\eqref{eq:ah-alg-graph-curvature} contributes
\(2p_Q(w,\mathrm dt)\) after division by two.  For the conformal
Laplacian term, the equality
\(\mathrm d\log(e^{2t}l)=(p+2)\mathrm dt+A\) yields
\begin{align*}
 -4\overline\Delta t-4|\mathrm dt|_{\overline g}^2
   ={}&-4u^{-1}\operatorname{div}_g(uA_d\nabla t)\\
     &+4(p+1)|\mathrm dt|_{A_d}^2
       +4\langle A,\mathrm dt\rangle_{A_d}.
\end{align*}
The cosmological terms are \(-3+3e^{4t}\).  Consequently, if
\(V_{\mathrm{pre}}=u(p_Qw+4A_d\nabla t)\), then
\begin{align}
 \frac12e^{4t}(R_{\widehat g}+6)
 ={}&8\pi\bigl(\mu+J(w)\bigr)+3(e^{4t}-1)
       +4\langle A,\mathrm dt\rangle_{A_d}\notag\\
 &+\frac12[Q,Q]+2p_Q(w,\mathrm dt)
       +4(p+1)|\mathrm dt|_{A_d}^2
       -u^{-1}\operatorname{div}_g V_{\mathrm{pre}}.
 \label{eq:ah-alg-before-trace}
\end{align}

We next incorporate the trace equation into the flux.  Direct
differentiation gives
\begin{equation}
 \nabla_iw_j=H^f_{ik}(A_d)_{kj}+dB_iw_j,\qquad
 \mathrm d\log\sqrt D=H^f(w,\cdot)+vB.
 \label{eq:ah-alg-w-derivatives}
\end{equation}
Since \(u=e^{2t}l\sqrt D\), these identities imply
\begin{equation}
 u^{-1}\operatorname{div}_g(uw)
   =F+(1-\chi)q_1-\operatorname{tr}_gK
                   +2(p+1)a x_t .
 \label{eq:ah-alg-uw-divergence}
\end{equation}
One way to see the cancellation of \(A\) is to first obtain
\[
 u^{-1}\operatorname{div}_g(uw)
   =\operatorname{tr}_g H^f+2B(w)+2\mathrm dt(w),
\]
and then use
\(\operatorname{tr}_gS
 =\operatorname{tr}_gK+\operatorname{tr}_gH^f+2A(w)
 =F+(1-\chi)q_1\).
Define \(V=V_{\mathrm{pre}}+uwF\).  The divergence of the added term contributes
\[
 w(F)+F\bigl[F+(1-\chi)q_1-\operatorname{tr}_gK
                         +2(p+1)a x_t\bigr]
\]
to Equation~\eqref{eq:ah-alg-before-trace}.

We verify the resulting quadratic expression explicitly.  Write
\(z=\operatorname{tr}_{e^\perp}P=F-\chi q_1\).  The three blocks
of \(Q\) are \(P\), \(M+pa y\), and \(q_1+2pa x_t\), so
\[
 \frac12[Q,Q]
   =\frac12|P^{\mathrm{tf}}|^2-\frac14z^2
      +|M+pa y|^2-zq_1-2pa x_tz ,
\]
and
\[
 2p_Q(w,\mathrm dt)
    =-2a x_tz+2a\langle M,y\rangle+2pv|y|^2 .
\]
Adding the terms in
Equation~\eqref{eq:ah-alg-before-trace} and
Equation~\eqref{eq:ah-alg-uw-divergence}, and setting aside
\(-F\operatorname{tr}_gK\), gives
\begin{align*}
 &\frac12|P^{\mathrm{tf}}|^2
   +|M+(p+1)a y|^2+[4(p+1)-v]|y|^2\\
 &\quad+4(p+1)d x_t^2+2(p+1)a\chi x_tq_1
       +\frac34F^2-\frac{\chi}{2}Fq_1
       +\left(\chi-\frac{\chi^2}{4}\right)q_1^2 .
\end{align*}
Completing the \(x_t\)-square leaves the coefficient
\[
 \chi-\frac{\chi^2}{4}
      -\frac{(p+1)v\chi^2}{4d}
   =\frac{12d}{(4+pv)^2}
\]
in front of \(q_1^2\).  Completing the \(F\)-square then gives
exactly Equation~\eqref{eq:ah-alg-T}.  Equivalently,
\(\mathcal T\) is the coordinate-free expression
\[
 \frac12[Q,Q]+2p_Q(w,\mathrm dt)
   +4(p+1)|\mathrm dt|_{A_d}^2
   +F\bigl[F+(1-\chi)q_1+2(p+1)a x_t\bigr].
\]
The apparently directional terms in this last expression are smooth:
\((1-\chi)q_1=((p+4)/(4+pv))S(w,w)\) and
\(a x_t=\mathrm dt(w)\).  This proves the asserted extension through
zero gradient.

Finally, Equation~\eqref{eq:ah-alg-w-derivatives} gives
\[
 4\,\mathrm dZ=(4+pv)\mathrm dt+H^f(w,\cdot)+vA.
\]
Using
\((1-\chi)(4+pv)=(p+4)v\), the preceding definition of \(V\)
reduces to
\begin{equation}
 \frac Vu=p_Qw+4A_d\nabla t+wF
       =A_\chi\bigl(4\nabla Z+K(w,\cdot)^\sharp+A(w)w\bigr).
 \label{eq:ah-alg-flux-reduction}
\end{equation}
This proves both the stated flux formula and
Equation~\eqref{eq:ah-alg-curvature-identity}.
\end{proof}

\begin{proposition}[The spatial primitive]
\label{prop:ah-alg-primitive}
Suppose the lapse is the profile defined in
Section~\ref{subsec:ah-profiles}; in particular
\(l_0(t,x)=e^{n(x)t}\) for \(t\le0\) and \(l_0(0,x)=1\).
Define
\begin{gather}
 G(t,x)=\lambda(x)\int_{-\infty}^t L(z,x)\,\mathrm dz,\qquad
 G_c(x)=G(0,x)=\frac{\lambda(x)}{n(x)+2},\notag\\
 V_*=V-4\sqrt D\,A_d\nabla_xG+4\nabla G_c .
 \label{eq:ah-alg-primitive-definitions}
\end{gather}
Assume the scalar equation of Equation~\eqref{eq:ah-system} is written
as
\[
 u^{-1}\operatorname{div}_gV
 =3(e^{4t}-1)+4\langle A,\mathrm dt\rangle_{A_d}+\mathcal P,
\]
where, for the physical system,
\[
 \mathcal P=d_0(\mathcal T+\tau A_Ra\sigma)
            +\rho-m_0 n^{P_0}(\rho_0+v),\qquad \sigma=|\mathrm df|.
\]
Put
\[
 b_1=\frac{G}{\lambda L},\qquad
 s_G=\mathrm d_x\log G,\qquad
 X_1=x_t+\frac{a q_1}{4+pv}.
\]
Then
\begin{equation}
 u^{-1}\operatorname{div}_gV_*=\mathcal P+\mathcal R,
 \label{eq:ah-alg-primitive-divergence}
\end{equation}
where
\begin{align}
 \mathcal R={}&3(e^{4t}-1)-4b_1E_G
           +\frac{4\sqrt d}{\lambda L}\Delta_gG_c\notag\\
 &+4b_1\left\{
       s_G(w)\bigl(F-\operatorname{tr}_{A_d}K+pdaX_1\bigr)
       -pv\langle(s_G)_{e^\perp},y\rangle
       \right\},
 \label{eq:ah-alg-primitive-residual}\\
 E_G={}&\frac{\Delta_xG-\nabla_x^2G(w,w)}{G}
              +v\langle s_G,A\rangle_{A_d}.
 \label{eq:ah-alg-EG}
\end{align}
Here \(\Delta_xG\) and \(\nabla_x^2G\) differentiate \(x\mapsto
G(t,x)\) with \(t\) fixed.  If \(\lambda\) and \(n\) are spatially
constant on a boundary collar, then \(V_*=V\) on that collar.
\end{proposition}

\begin{proof}
For a fixed spatial function \(\phi\), write
\[
 \mathcal L\phi
 =D^{-1/2}\operatorname{div}_g(\sqrt D\,A_d\nabla\phi).
\]
Differentiating \(A_d=I-w\otimes w\) with
Equation~\eqref{eq:ah-alg-w-derivatives} gives
\[
 \mathcal L\phi
    =\operatorname{tr}_{A_d}\nabla^2\phi
        +\mathfrak q(\phi),
\]
where the drift vector is
\begin{equation}
 \mathfrak q
 =vA_d(A+p\,\mathrm dt)^\sharp
  -w\left(F+(d-\chi)q_1-\operatorname{tr}_{A_d}K
                         +2dpa x_t\right).
 \label{eq:ah-alg-primitive-drift}
\end{equation}
To check every contraction, the coefficient of \(\nabla\phi\) before
substituting \(S\) is
\[
 A_d\nabla\log\sqrt D
       -(\operatorname{div}_gw)w-\nabla_ww
 =vA_dB^\sharp
       -w\bigl(\operatorname{tr}_{A_d}H^f+2dB(w)\bigr).
\]
Now
\[
 \operatorname{tr}_{A_d}H^f
  =F+(d-\chi)q_1-\operatorname{tr}_{A_d}K-2dA(w),
\]
which proves Equation~\eqref{eq:ah-alg-primitive-drift}.

In applying this computation to \(\nabla_xG(t(x),x)\), there is one
additional differentiation in \(t\).  The defining integral gives
\[
 G_t=\lambda L,\qquad
 \partial_t\nabla_xG=\nabla_x(\lambda L)=\lambda L A^\sharp.
\]
It follows that
\begin{align*}
 \frac1u\operatorname{div}_g(\sqrt D\,A_d\nabla_xG)
  ={}&b_1\left\{
        \frac{\Delta_xG-\nabla_x^2G(w,w)}G
                 +\langle\mathfrak q,s_G\rangle
        \right\}
      +\langle A,\mathrm dt\rangle_{A_d}.
\end{align*}
Thus the factor \(-4\) in the primitive correction cancels exactly
the prescribed \(4\langle A,\mathrm dt\rangle_{A_d}\).
The \(A\)-part of \(\mathfrak q\) produces \(E_G\).  The remaining
\(\mathrm dt\)-dependent terms are
\[
 4b_1\left\{
   s_G(w)\bigl[(d-\chi)q_1+2dpa x_t\bigr]
                  -pv\langle s_G,\mathrm dt\rangle_{A_d}
       \right\}.
\]
Use
\[
 d-\chi=\frac{pdv}{4+pv},\qquad
 s_G(w)=a(s_G)_e,\qquad
 \langle s_G,\mathrm dt\rangle_{A_d}
       =\langle(s_G)_{e^\perp},y\rangle+d(s_G)_e x_t.
\]
These terms become
\[
 4b_1\left\{s_G(w)pda
                \left(x_t+\frac{a q_1}{4+pv}\right)
                    -pv\langle(s_G)_{e^\perp},y\rangle\right\},
\]
as asserted.  The last correction in \(V_*\) contributes
\(4\Delta_gG_c/u=4\sqrt d\,\Delta_gG_c/(\lambda L)\).
This establishes
Equations~\eqref{eq:ah-alg-primitive-divergence}--\eqref{eq:ah-alg-EG}.
On a collar with spatially constant \(\lambda,n\), both
\(\nabla_xG\) and \(\nabla G_c\) vanish.  The boundary assertion
follows.
\end{proof}

\begin{lemma}[A coercive part of the quadratic form]
\label{lem:ah-alg-coercivity}
For \(p\ge0\) and \(0<d\le1\), the form in
Equation~\eqref{eq:ah-alg-T} satisfies
\begin{equation}
 \mathcal T\ge\frac23F^2+[4(p+1)-v]|y|^2
                    +4(p+1)dX_1^2 .
 \label{eq:ah-alg-coercivity}
\end{equation}
\end{lemma}

\begin{proof}
The first two terms in Equation~\eqref{eq:ah-alg-T} are nonnegative.
The last two terms, keeping the \(X_1\)-square intact, combine to
\[
 \frac34F^2-\frac{2d}{4+pv}Fq_1
                +\frac{12d}{(4+pv)^2}q_1^2.
\]
Completing their square gives the exact identity
\begin{equation}
 \frac34F^2-\frac{2d}{4+pv}Fq_1
       +\frac{12d}{(4+pv)^2}q_1^2
   =\frac{9-d}{12}F^2
      +\frac{12d}{(4+pv)^2}
          \left(q_1-\frac{4+pv}{12}F\right)^2.
 \label{eq:ah-alg-q-minimum}
\end{equation}
Since \((9-d)/12\ge2/3\), the result follows.
\end{proof}

\paragraph{Asymptotic hypotheses for the flux calculation.}
For a tensor \(T\), the notation \(T=O_j(r^{-\gamma})\) means that
its \(b\)-norm and the \(b\)-norms of its covariant derivatives through
order \(j\) are \(O(r^{-\gamma})\).  The following statement concerns
each fixed choice of the profile parameters.  Its constants need not
be uniform as those parameters vary.

\begin{proposition}[Passage to the four metric fluxes]
\label{prop:ah-alg-mass-passage}
Suppose \(g\) has an asymptotically hyperbolic end with background \(b\)
and four finite metric fluxes \(p_a(g)\), \(a=0,1,2,3\), in the
designated chart, as defined in Equation~\eqref{eq:setup-flux}.
Assume that, outside a compact set, \(n\) is a
constant and
\(\lambda=\sqrt{A_0^2+r^2}\), with fixed \(A_0>0\).  Assume the
differentiated end bounds
\begin{equation}
 \begin{gathered}
 g-b=O_2(r^{-\kappa}),\qquad K=O_1(r^{-\kappa}),\\
 t=O_2(r^{-s}),\qquad w=O_2(r^{-\beta}),\\
 \kappa>\frac32,\qquad s>\frac32,\qquad\beta>\frac32,\qquad
 \kappa+\beta>3 .
 \end{gathered}
 \label{eq:ah-alg-end-W}
\end{equation}
Finally, suppose that for some \(\delta>0\),
\begin{equation}
 R_{\widehat g}+6=O(r^{-3-\delta}).
 \label{eq:ah-alg-end-R}
\end{equation}
Then all four metric fluxes of \(\widehat g\) exist and are finite.
For each \(a=0,1,2,3\),
\begin{equation}
 p_a(\widehat g)-p_a(g)
   =-\frac1{2\pi}
       \lim_{r\to\infty}
       \int_{S_r}(V_a\,\mathrm dt-t\,\mathrm dV_a)(\nu_b)
                         \,\mathrm dA_b .
 \label{eq:ah-alg-four-flux}
\end{equation}
Moreover,
\begin{equation}
 V_*=4(\lambda\nabla_gt-t\nabla_g\lambda)+o(r^{-2}),
 \label{eq:ah-alg-Vstar-asymptote}
\end{equation}
and the time-component identity is
\begin{equation}
 p_0(\widehat g)-p_0(g)
   =-\frac1{8\pi}\lim_{r\to\infty}
        \int_{S_r}g(V_*,\nu_g)\,\mathrm dA_g .
 \label{eq:ah-alg-time-flux}
\end{equation}
\end{proposition}

\begin{proof}
We first show that the nonlinear terms have no mass flux.  Since
\(D=(1-|w|_g^2)^{-1}\), the graph correction is exactly
\[
 l^2\mathrm df\otimes\mathrm df
        =D\,w^\flat\otimes w^\flat .
\]
Consequently, with differentiated bounds through order one,
\begin{equation}
 \widehat g-g=4t\,b+\mathcal E,\qquad
 \mathcal E
    =O_1(r^{-2s})+O_1(r^{-s-\kappa})+O_1(r^{-2\beta}).
 \label{eq:ah-alg-metric-error}
\end{equation}
For the four background static potentials,
\[
 |\nabla_b^jV_a|_b=O(r)\quad(j=0,1,2),\qquad
 |S_r|_b=4\pi r^2.
\]
Thus a metric error \(O_1(r^{-j})\) contributes at most
\(O(r^{3-j})\) to its flux.  The assumptions in
Equation~\eqref{eq:ah-alg-end-W} make all three contributions in
Equation~\eqref{eq:ah-alg-metric-error} tend to zero.

Write \(\mathbb U_b(V,e)\) for the one-form
\[
 V(\operatorname{div}_b e-\mathrm d\operatorname{tr}_b e)
       +(\operatorname{tr}_b e)\mathrm dV
       -e(\nabla_bV,\cdot)
\]
appearing in the definition of the metric mass.  In three dimensions,
\begin{equation}
 \mathbb U_b(V,4t\,b)=-8(V\,\mathrm dt-t\,\mathrm dV).
 \label{eq:ah-alg-conformal-flux}
\end{equation}
Indeed \(\operatorname{div}_b(4tb)=4\mathrm dt\),
\(\operatorname{tr}_b(4tb)=12t\), and
\((4tb)(\nabla_bV,\cdot)=4t\,\mathrm dV\).  Combining
Equation~\eqref{eq:ah-alg-conformal-flux} with
Equation~\eqref{eq:ah-alg-metric-error} proves the asserted
comparison of the finite-radius flux integrals.

We establish existence of the limiting fluxes before taking their
difference.  Put \(\widehat e=\widehat g-b\) and
\(\nu=\min(\kappa,s,2\beta)>3/2\).  Then
\(\widehat e=O_2(r^{-\nu})\).  The scalar-curvature linearization at
the curvature \(-1\) background is
\[
 D R_b(e)=\operatorname{div}_b\operatorname{div}_b e
           -\Delta_b\operatorname{tr}_b e+2\operatorname{tr}_b e.
\]
Since \(\nabla_b^2V_a=V_a b\) and \(\Delta_bV_a=3V_a\),
differentiating the defining one-form gives exactly
\begin{equation}
 \operatorname{div}_b\mathbb U_b(V_a,\widehat e)
                   =V_a D R_b(\widehat e).
 \label{eq:ah-alg-linearized-divergence}
\end{equation}
The Taylor remainder in scalar curvature has the estimate
\[
 \bigl|D R_b(\widehat e)-(R_{\widehat g}+6)\bigr|
 \le C\bigl(
       |\widehat e|\,|\nabla_b^2\widehat e|
       +|\nabla_b\widehat e|^2+|\widehat e|^2
       \bigr)
 =O(r^{-2\nu}).
\]
The constant is uniform on a sufficiently distant tail, where
\(\widehat g\) and \(b\) are uniformly equivalent.  As
\[
 \mathrm dV_b=\frac{r^2}{\sqrt{1+r^2}}\,\mathrm dr\,\mathrm d\omega,
\]
the scalar-curvature contribution and the Taylor remainder in the
right side of Equation~\eqref{eq:ah-alg-linearized-divergence} are
absolutely integrable there:
\[
 \int^\infty r^{-1-\delta}\,\mathrm dr<\infty,\qquad
 \int^\infty r^{2-2\nu}\,\mathrm dr<\infty .
\]
The divergence theorem on annuli therefore makes each family of
flux integrals a Cauchy family.  All four limits exist and are finite.
Taking the difference from the given fluxes of \(g\) now proves
Equation~\eqref{eq:ah-alg-four-flux}.

It remains to identify the corrected vector field with the time
potential.  On the ultimate tail, set
\[
 H_n(t)=\int_{-\infty}^t L(z)\,\mathrm dz.
\]
The constant \(n\) is fixed here, and
\[
 H_n(0)=\frac1{n+2},\qquad H_n'(0)=1,\qquad
 \nabla_xG=H_n(t)\nabla_g\lambda .
\]
Taylor expansion at zero gives
\[
 \nabla_xG
   =\left(\frac1{n+2}+t+O(t^2)\right)\nabla_g\lambda,\qquad
 \nabla_gG_c=\frac{\nabla_g\lambda}{n+2}.
\]
On the same tail \(A=\mathrm d\log\lambda\), so
\[
 \sqrt D A_d=I+O(r^{-2\beta}),\qquad
 \frac u\lambda=1+O(r^{-s})+O(r^{-2\beta}),\qquad
 \nabla Z=\nabla t+O(r^{-2\beta}).
\]
Substituting these bounds into
Equation~\eqref{eq:ah-alg-V} and
Equation~\eqref{eq:ah-alg-primitive-definitions} yields the explicit
error ledger
\begin{equation}
 V_*-4(\lambda\nabla_gt-t\nabla_g\lambda)
   =O(r^{1-2s})+O(r^{1-2\beta})
                         +O(r^{1-\kappa-\beta}).
 \label{eq:ah-alg-vector-error}
\end{equation}
All three terms are \(o(r^{-2})\) by
Equation~\eqref{eq:ah-alg-end-W}.  In particular their integrated
normal fluxes tend to zero.

Finally,
\[
 \lambda-V_0=O_1(r^{-1}),\qquad V_0=\sqrt{1+r^2}.
\]
Replacing \(V_0\) by \(\lambda\) in the conformal flux integral costs
\(O(r^{1-s})=o(1)\).  Replacing \(\nu_b,\mathrm dA_b\), and the
background gradient by their \(g\)-counterparts costs
\(O(r^{3-s-\kappa})=o(1)\), since \(s+\kappa>3\).
Equation~\eqref{eq:ah-alg-time-flux} follows from
Equations~\eqref{eq:ah-alg-four-flux} and
\eqref{eq:ah-alg-Vstar-asymptote}.  The proof treats the four static
potentials separately; it does not assume a causal character for the
mass covector of \(\widehat g\).
\end{proof}

\begin{lemma}[Scalar-curvature decay from the end equation]
\label{lem:ah-alg-scalar-decay}
Suppose the differentiated bounds of
Equation~\eqref{eq:ah-alg-end-W} hold,
\(R_g+6=O(r^{-3-\delta})\), and
\[
 2s>3+\delta,\qquad 2\beta>3+\delta.
\]
Suppose also that the scalar end equation has the form
\begin{equation}
 \Delta_gZ+C(Z)|\mathrm dZ|_g^2
      =\frac34(e^{4Z}-1)+O(r^{-3-\delta}),
 \label{eq:ah-alg-end-equation}
\end{equation}
where \(C\) is smooth and bounded near zero for the fixed profiles.
Then Equation~\eqref{eq:ah-alg-end-R} holds.  In the deformation
equations one takes
\[
 C(Z)=p(Z)+2-d_0(Z)\bigl(p(Z)+1\bigr)
\]
on the tail with spatially constant \(n\).
\end{lemma}

\begin{proof}
Since \(D=(1-|w|^2)^{-1}\), we have
\(Z-t=O_2(r^{-2\beta})\).  The assumed end equation therefore gives
\[
 \Delta_gt+C(t)|\mathrm dt|_g^2
      =\frac34(e^{4t}-1)+O(r^{-3-\delta}).
\]
First apply the conformal scalar-curvature formula to \(e^{4t}g\):
\begin{align*}
 e^{4t}\bigl(R_{e^{4t}g}+6\bigr)
   &=R_g+6-8\Delta_gt-8|\mathrm dt|_g^2+6(e^{4t}-1)\\
   &=R_g+6+8\bigl(C(t)-1\bigr)|\mathrm dt|_g^2
                                      +O(r^{-3-\delta}).
\end{align*}
The gradient term is \(O(r^{-2s})=O(r^{-3-\delta})\).
The difference
\(\widehat g-e^{4t}g=e^{4t}D\,w^\flat\otimes w^\flat\)
is \(O_2(r^{-2\beta})\).  The local scalar-curvature formula,
on the uniformly equivalent end metrics with bounded curvature,
then gives
\(R_{\widehat g}-R_{e^{4t}g}=O(r^{-2\beta})\).
This proves the required decay.
\end{proof}

\begin{remark}
The scalar-curvature hypothesis in the mass proposition is essential
to the convergence argument.  Bounds \(t=O_2(r^{-s})\) for every
\(s<3\), by themselves, would allow \(t=r^{-3}\log r\), whose
conformal time-flux integral diverges logarithmically.  In the
construction, Lemma~\ref{lem:ah-alg-scalar-decay} supplies the missing
integrability from the scalar equation.  The estimates of the
primitive residual are then used separately to determine the sign of
the limiting flux in Equation~\eqref{eq:ah-alg-time-flux}.
\end{remark}


\subsection{A specified continuation and uniform height control}
\label{subsec:ah-homotopy}

Choose \(j\in C^\infty(\mathbb R;[0,1])\) equal to zero on
\(t\le0\) and one on \(t\ge1\), and put
\[
 \mathcal D=A_\chi\{K(w,\cdot)+A(w)w\},\qquad
 V_\vartheta=u(4A_\chi\nabla Z+\vartheta\mathcal D),
 \quad 0\le\vartheta\le1.
\]
Before the final scalar stage the inner condition is
\begin{equation}\label{eq:ah-homotopy-boundary}
 \frac{(V_\vartheta)_\nu}{u}
 =[1-(1-\vartheta)j(t)]
       [\,\mathcal B-(1-\vartheta)\mathcal D_\nu\,].
\end{equation}
The outer values are always \(f=Z=0\).  The scalar source is always
the expression \(\Xi\) in Equation~\eqref{eq:ah-system}, evaluated
using the current trace \(F\).

Here are explicit trace corrections.  On a black collar take a
negative multiple of a spatial cutoff, a directional cutoff
supported on \(w\cdot\nu_s<-1/2\), and a nonincreasing height cutoff
equal to one near \(b_-\) and zero near \(b_+\).  On a white collar
use a positive multiple, a directional cutoff supported on
\(w\cdot\nu_s>1/2\), and a nondecreasing height cutoff supported
near \(b_+\).  Their sum is denoted by \(D_0(x,w,h)\).
It is nondecreasing in \(h\), vanishes in compatible directions,
and is nonpositive at sufficiently low heights and nonnegative at
sufficiently high heights.  Choose its amplitudes larger than
\(2\sup_{\partial\Omega}|H|+1\).  The limiting directional
inequalities at the assigned face height then are
\begin{equation}\label{eq:ah-directional-barriers}
 F(x,+\nu,b)\ge P_B+H,\qquad
 F(x,-\nu,b)\le P_B-H .
\end{equation}
These are evaluations at \(|w|=1\), hence \(d=\chi=0\).
Let \(b(x)\) be a compactly supported extension of the face values,
constant on smaller collars, and let
\(D_1(x,w)=M_1w\cdot\nu_s\) on those collars, cut off smoothly,
where \(M_1>\sup|H|+1\).

The four stages are:
\begin{enumerate}[label=(\roman*)]
 \item Decrease \(\vartheta\) from one to zero, retaining \(F=h+C\)
       and the original face heights.
 \item Keep \(\vartheta=0\) and the face heights, and set
       \(F=h+C+\alpha D_0(x,w,h)\), with \(\alpha\) increasing from
       zero to one.
 \item Keep \(\vartheta=0\), prescribe \(h|_{\partial\Omega}=(1-\zeta)b\),
       and set
       \begin{equation}\label{eq:ah-third-stage}
       F=h+(1-\zeta)\{C+D_0(x,w,h+\zeta b(x))\}
          +\zeta\{\tr_{A_\chi}K+D_1(x,w)\},
       \qquad 0\le\zeta\le1.
       \end{equation}
 \item Keep the trace problem at \(\zeta=1\), and multiply both
       \(\Xi\) and the right side of
       Equation~\eqref{eq:ah-homotopy-boundary} by a scalar decreasing
       from one to zero.
\end{enumerate}
The prescriptions agree at their junctions.  At fixed \(x,Z,df\),
each trace right side has derivative at least one in \(h\).
The inequalities in Equation~\eqref{eq:ah-directional-barriers}
persist in the third stage by convex combination.  In the last
stage \(f=0\): at an extremum \(w=0\), \(D_1=0\), and the trace
equation has the strict height sign.  Hence \(t=Z\) there.

\begin{lemma}[Height, confinement, and the order of compact choices]
\label{lem:ah-heights}
For all sufficiently large \(N\), and all sufficiently distant outer
truncations at fixed profiles, every smooth solution in the first
three stages satisfying \(t\ge-\epsilon\) obeys
\begin{equation}\label{eq:ah-height-bounds}
 |\tau f|\le B,\qquad |h|\le Br^{-\beta}\quad(r\ge r_0).
\end{equation}
The outer normal slope obeys the differentiated rate of the same
radial barrier, and \(t>-\epsilon\) on the outer face.  The constants
\(B,r_0\) may be chosen independently of collar widths and collar
derivatives, uniformly while the thresholds and correction amplitudes
range in fixed bounded sets.

The thresholds and offsets can consequently be fixed so that, on
every physical solution obtained below,
\begin{equation}\label{eq:ah-offset-budget}
 |F\tr K|+|\nabla C|\le2\rho .
\end{equation}
The fixed region on which \(n=N\) can be chosen to contain every
compact sublevel set used in the boundary slope argument.
\end{lemma}

\begin{proof}
At an interior extremum of \(f\), \(w=0\), so \(A\), \(D_0\), and
\(D_1\) contribute nothing to the trace equation.  In stage three
it reads
\[
 \tr(l\Hess f)=\tau A_Rf+(1-\zeta)(C-\tr K).
\]
In the first two stages the same formula has \(\zeta=0\).
Since \(A_R\ge1\), the maximum principle and the boundary values
give
\[
 |\tau f|\le
 \max\{|b_-|,|b_+|,\sup|\tr K|+\sup|C|\}.
\]
No derivative of a collar correction has been used.  The necessary
correction amplitudes are uniformly bounded: on a black face
\(H=c_b-\tr_TK\), and the white face has the analogous bound.
All supports can be kept in a fixed compact enlargement by
shortening the individual collars.

Take \(r_0\) outside that enlargement and the fixed lapse transition.
Compare \(f\) with
\[
 \pm\frac{B_1}{\tau}
 \left\{\frac{r^{-\beta}}{A_R(r)}
       -\frac{T^{-\beta}}{A_R(T)}\right\}
\]
on a truncation at \(r=T\).  At a contact the axis is radial.
Put \(k=\beta+\partial_{\log r}\log A_R\), so
\(\beta\le k\le\beta+1<4\) and its derivative is nonnegative.
The Hessian term, together with the radial part \(2\chi A(w)\)
coming from \(s=d\log\lambda\), has upper bound
\[
 a\{-2+\chi(k-2)+o(1)\}
\]
for the positive decreasing test.  This is bounded above by a
fixed negative multiple of \(a\), uniformly for \(0\le\chi\le1\).
The derivative of k has favorable sign.

The only additional spatial-profile term occurs where n varies.
Positive t helps, because \(pt\,d\log n\) points outward.
For negative t put \(X_0=-nt\).  The adverse term is at most
\(C\chi aX_0\).  If
\(X_0>C_1\log(r/\tau)\), then \(l_0=e^{-X_0}\) and the tested
gradient gives
\[
 \chi aX_0\le
 C X_0e^{-X_0}\frac{\lambda}{\tau A_R}r^{-\beta}
 =o(r^{-\beta})
\]
by taking the fixed \(C_1\) large.  In the complementary regime
\(\chi a\le C/\sqrt n\), so
\[
 \chi aX_0\le
 C\frac{\log r+\log N}{\sqrt n}=o(1)r^{-\beta},
\]
uniformly on the varying region because \(q>2\beta\).
There the shift at \(T\) is negligible after \(T\) is sufficiently
large at fixed profiles.

The tensor term is \(O(r^{-\kappa})\), dominated by the positive
test height since \(\beta<\kappa\), except near \(T\).
Near \(T>2R\), n is constant, and on the test
\(a\ge c_Nr^{-\beta}\) whenever the tested slope is small:
\(\lambda/A_R\) and \(l_0\) have positive lower bounds at fixed
profiles.  Thus the negative Hessian contribution dominates the
tensor term there as well.  For larger slope it is stronger.
The negative test gives the reverse inequality.  Choose \(B_1\)
to dominate the already proved bound on the fixed sphere \(r_0\);
this uses no collar derivatives.  Height monotonicity proves
Equation~\eqref{eq:ah-height-bounds}.  At the outer face the barrier
derivative tends to zero in the quantity \(l|df|\), at fixed
profiles.  The relation \(Z=t+\log(1+l^2|df|^2)/8=0\) then excludes
\(t=-\epsilon\), once \(T\) is sufficiently large.

We now specify the compact choices without circular dependence.
Use this common B to choose a compact set Q outside which
\(Br^{-\beta}|\tr K|\le\rho\).
This is possible because \(\beta+\kappa>3+\delta\).
Choose the two negative thresholds sufficiently close to zero,
then their collars, and finally sufficiently small offset amplitudes,
so that
\[
 |(h+C)\tr K|+|\nabla C|\le\rho
 \quad\hbox{on Q whenever } b_-\le h\le b_+ .
\]
The constants controlling derivatives of the now fixed collars
need not be uniform in this choice.  The next lemma proves precisely
the required physical height interval on Q.  Outside Q the end
bound proves Equation~\eqref{eq:ah-offset-budget}, after making
the support of C lie in Q.  Finally enlarge the fixed constant-n
region so that Equation~\eqref{eq:ah-height-bounds} rules out
heights near \(b_-\) and \(b_+\) beyond it.  This radius is fixed
before N and R.
\end{proof}

\begin{lemma}[Compact separation and polynomial collar estimates]
\label{lem:ah-collars}
For the fixed choices above, all sufficiently large \(N\), and smooth
solutions satisfying the floor condition \(t\ge-\epsilon\), the
first two stages have constants \(c,C>0\), independent of \(N\)
and outer truncation, such that
\[
 \frac c\tau\le\partial_\nu f\le\frac C\tau
 \quad\hbox{on black faces},\qquad
 \frac c\tau\le-\partial_\nu f\le\frac C\tau
 \quad\hbox{on white faces}.
\]
In stage three, \(|\partial_\nu f|\le C/\tau\).
On each fixed compact set needed in
Equation~\eqref{eq:ah-offset-budget}, the physical heights satisfy
\(b_-\le h\le b_+\).
On the fixed collars \(\mathcal C\), for \(\phi\ge0\),
\begin{equation}\label{eq:ah-polynomial-trace}
 \begin{split}
 \int_{\partial\Omega}\lambda L\phi
 &\le CN^m\int_{\mathcal C}
 u\{(1+\sqrt{\mathcal T})\phi+|d\phi|_{A_\chi}\},\\
 \int_{\partial\Omega}\phi
 &\le CN^m\int_{\mathcal C}
 \sqrt D\{(1+\sqrt{\mathcal T})\phi+|d\phi|_{A_\chi}\}.
 \end{split}
\end{equation}
The finite exponent m is independent of \(B_0,R,n_\infty\);
fixed collar geometry affects the constant C.
\end{lemma}

\begin{proof}
We verify the extension of Lemma~\ref{lem:fc-compact-separation}
rather than assume an asymptotically flat hypothesis.  A putative
failure supplies \(N_i\to\infty\) and outer radii tending to infinity.
On the relevant compact range \(A_R=1,n=N\), and
\(l/\tau\ge c\ell/\tau=cN^{B_0/4}\to\infty\).
At a fixed nonzero-gradient lower test \(\phi\) of a lower relaxed
height limit,
\[
 d_i,\chi_i\longrightarrow0,\qquad a_i\longrightarrow1,\qquad
 \frac{l_i}{\tau_i\sqrt{D_i}}\longrightarrow|d\phi|^{-1},
 \qquad 2\chi_iA(w_i)\longrightarrow0.
\]
The last term is the only new trace term on this compact range,
where \(A=s\) is bounded.  The limiting inequality is therefore
\[
 \frac{\tr_{(d\phi)^\perp}\Hess\phi}{|d\phi|}
       +\tr_{(d\phi)^\perp}K\le k'+C_b
\]
at heights below \(k'\), with \(b_-<k'<0\) close to \(b_-\).
The low-height \(D_0\) term is nonpositive.  The lower relaxed passage,
increasing logistic reparametrizations, and
Lemma~\ref{lem:fl-support-enclosure} now apply exactly under the
hypotheses established in Lemma~\ref{lem:fc-compact-separation}.
In particular each height is extended constantly as \(b_+\) across
a white face before taking the relaxed limit.  Low test contacts
are therefore interior, even if a boundary layer is present.
The reversed white face has expansion \(-c_w>0\), and excludes
contact of the resulting low closed sublevel set.  The end bound
makes that set compact.  After adjoining the full black region,
right continuity of the black threshold family excludes its meeting
any compact set away from that region.  Reversing signs proves the
upper separation, including the empty-white-region case.

On an actual boundary collar \(A=0\), \(n=N\), \(A_R=1\), and
\(\lambda=\lambda_0\) is constant.  For a test
\(h=b_-+\psi(s)\) of fixed positive slope,
\[
 d\le C(\tau/\ell)^2=C N^{-B_0/2},
 \qquad N\chi\ge c d.
\]
The trace residual is the leaf expansion plus \(O(d)\) and the
term \(a\chi|ds|(\log\psi')'\).
Use
\[
 \psi'=m_1\exp\!\left(\int_0^s\gamma_N\right)
 \quad\hbox{or}\quad
 \psi'=M_2\exp\!\left(-\int_0^s\gamma_N\right),
\]
where \(\gamma_N\) is a fixed large multiple of N on
\([0,1/N]\), vanishes beyond \(2/N\), and has bounded integral.
The first term dominates \(O(d)\) on the short interval.
Outside it the offset margin is a fixed multiple of s and
dominates d when \(B_0>2\).  Compact separation orders the lower
barrier on the far end of the collar; fixed large slopes order
the upper barrier.  The third-stage directional inequalities
give the same upper slope barriers of both signs.
This verifies all hypotheses and all changed small parameters in
the proof of Lemma~\ref{lem:fc-collar-slopes}; it does not require
the special flat choice \(\tau=\ell^{3/2}\).

For the trace inequalities use the signed field
\(\lambda_0Lw\), whose boundary flux has magnitude at least
\(\lambda_0L/2\).  Direct differentiation on these collars gives
\[
 \frac{\operatorname{div}(\lambda_0Lw)}u
 =\sqrt d\{\tr_{A_d}H^f+(dp+p+2)a x_t\}.
\]
With the L weight removed, the same formula has only \(dp\,a x_t\)
in its gradient term.  Coercivity of
Equation~\eqref{eq:ah-T}, equivalently
Lemma~\ref{lem:fl-alg-coercivity}, bounds these expressions by
\(CN^m(1+\sqrt{\mathcal T})\).
The gradient action is bounded by \(|d\phi|_{A_d}\), and
\(A_d\le(1+N/4)A_\chi\).
Integration with fixed collar cutoffs proves
Equation~\eqref{eq:ah-polynomial-trace}.
No inverse power of \(\ell\) occurs in this calculation.
\end{proof}

\subsection{Exclusion of the floor}
\label{subsec:ah-floor}

\begin{lemma}[Floor separation]\label{lem:ah-floor}
A fixed sufficiently small \(\delta_0>0\) and then a fixed finite
integer \(P_0\) can be chosen so that, for all sufficiently large
N and all sufficiently distant truncations at fixed profiles,
no smooth solution anywhere in the continuation can satisfy
\(t\ge-\epsilon\) and attain \(t=-\epsilon\).
Neither choice requires a fixed-profile Schauder constant.
\end{lemma}

\begin{proof}
At an interior contact \(dt=0,\Hess t\ge0\).  The Gauss equation of
the graph in the Riemannian warped product \(g+l^2ds_0^2\) gives
\begin{equation}\label{eq:ah-floor-gauss}
 \frac12e^{4t}R_{\hat g}
 \le \frac12R_g
   -v\left(\Ric(e,e)+\frac{\tr_\perp\Hess_xl}{l}\right)
   +\mathcal S\left(S-K+d^{-1}A(w)w\otimes w\right),
\end{equation}
where
\[
 \mathcal S(B)=\frac14(\tr_\perp B)^2
   -\frac12|B_{\perp\perp}^{\rm tf}|^2
   +dB_{ee}\tr_\perp B-d|B_{e\perp}|^2.
\]
To obtain this formula, the spatial lapse derivatives contribute
\(w\otimes A+A\otimes w\) to the graph second form, and its
vertical connection contributes \(d^{-1}A(w)w\otimes w\).
The remaining \(p\Hess t\) lapse curvature and conformal Hessian
terms are nonpositive at a contact.  The rank-one term changes
\(\mathcal S\) by only \(vA(w)\tr_\perp(S-K)\).

We give the derivative ledger that controls the additional terms.
At a floor contact, \(p=n\), \(d_xp=dn\), and
\[
 A=s-\epsilon n\,d\log n,\qquad
 \nabla w=H^fA_d+dA\otimes w,\qquad
 d\log u=H^f(w,\cdot)+(1+v)A .
\]
Positive-order derivatives of \(\log n\) are bounded.  Thus
\(\Hess_xl/l\), \(A\), and the needed spatial derivatives have
fixed-degree polynomial bounds in \(n\), uniformly in position.
The factor \(\epsilon<1\) introduces no larger growth.  Moreover
\[
 \partial_p\frac{4+p}{4+pv}
   =\frac{4d}{(4+pv)^2},\qquad
 dv=2d\{H^f(w,\cdot)+vA\}
\]
at contact.  Therefore the derivatives of \(A_\chi\) always carry
a \(d\) factor in the uncontrolled axial Hessian direction.
In \(d\log u\), that direction contracts with \(\chi\) in the flux.
The same is true for the new drift
\(A_\chi(A(w)w)=\chi A(w)w\).
A spatial derivative of \(A\) in this drift is contracted axially
and retains \(\chi\).

The algebraic floor comparison of
Lemma~\ref{lem:fl-alg-floor-coercivity}, applied to \(S\), and
Equation~\eqref{eq:ah-floor-gauss} therefore give, with universal
\(c>0\), small prescribed \(\eta>0\), and fixed finite powers \(m\),
the bounds
\[
 \begin{split}
 \mathcal T-\mathcal S
       \left(S-K+d^{-1}A(w)w\otimes w\right)
 &\ge c\mathcal T-Cn^m(\rho_0+v),\\
 w(F)&\ge\tau A_Ra\sigma-\eta\mathcal T-C_\eta n^m(\rho_0+v),\\
 u^{-1}|\operatorname{div}(u\mathcal D)|
 &\le\eta\mathcal T+C_\eta n^m(\rho_0+v).
 \end{split}
\]
Here \(|K|^2+|\nabla K|^2\lesssim\rho_0\), and the height
estimate bounds the squared end height by the same weight.
For example \(a|\nabla K|\) is bounded by \(v+|\nabla K|^2\).
The derivative of \(h=\tau A_Rf\) gives the displayed positive
term and \(h\,w(d\log A_R)\), whose absolute value is at most
\(Ca|h|\le C(v+\rho_0)\) for this floor estimate.
Derivatives of \(D_0,D_1,b,C\) are fixed compact coefficients;
the derivative in h is nonnegative and at least one.  These
observations enumerate all new spatial terms beyond the flat
floor calculation, including \(d_xp\), which is not zero here.

Combining these bounds with
Equation~\eqref{eq:ah-curvature} and the constraint gives
\begin{equation}\label{eq:ah-floor-divergence}
 u^{-1}\operatorname{div}V_\vartheta
 \ge 2\delta_0\mathcal T+\tau A_Ra\sigma
                  -Cn^m(\rho_0+v)
\end{equation}
for fixed sufficiently small \(\delta_0\).
The cosmological constants cancel in deriving this lower bound:
\(8\pi\mu-R_g/2-3=Q_g(K)/2\).
In the prescribed source the remaining explicit term
\(3(e^{4t}-1)\) is negative at the floor, where \(d_0=\delta_0\)
and \(m_0=1\).  Choose \(P_0>m+1\) and N sufficiently large
that its penalty dominates the error and \(2\rho+\rho_0\).
Equation~\eqref{eq:ah-floor-divergence} then contradicts the
scalar equation.

Outer contacts were excluded in Lemma~\ref{lem:ah-heights}.
At an inner contact \(j=0\), so
Equation~\eqref{eq:ah-homotopy-boundary} restores the full flux
condition.  Since \(A=0\) on the collar, its face identity is
\[
 d(H+4\partial_\nu t)=-N_Bd.
\]
It forces \(\partial_\nu t<0\), contrary to a minimum.
In the final scalar stage \(f=0,t=Z\).  At a floor minimum the
unscaled source is negative after the same penalty is enlarged;
at positive scale the inner derivative is negative as well.
At zero scale the homogeneous mixed problem gives \(Z=0\) by
multiplication by Z.  This proves separation for all stages.
\end{proof}

\subsection{Bounds, a compact equation, and exhaustion}
\label{subsec:ah-existence}

\begin{proposition}[Existence at fixed profiles]\label{prop:ah-existence}
Choose \(\delta_0,P_0\) as in Lemma~\ref{lem:ah-floor}.
There is a fixed sufficiently large \(B_0\), followed by a fixed
exponent \(a\) satisfying Equation~\eqref{eq:ah-R-conditions}, such
that for every sufficiently large \(N\) the physical system has a
smooth solution on the infinite exterior with \(t\ge-\epsilon\).
It is a smooth compact limit of solutions on sufficiently distant
spherical truncations. At fixed profiles, \(Z,t,D\), and all their
derivatives on hyperbolic unit patches are bounded; derivatives
of \(f\) are bounded after multiplication by the local value of
\(\lambda\). These bounds are uniform in outer truncation.
Moreover \(R_{\widehat g}\ge-6\), and its inner boundary has strictly
negative mean curvature.
\end{proposition}

\begin{proof}
We first prove estimates for all smooth floor-admissible solutions
of the continuation, without assuming their existence.

\paragraph{The principal divergence at high levels.}
Put \(B=A+p\,dt\). Differentiation gives
\[
 \nabla w=H^fA_d+dB\otimes w,\qquad
 d\log u=(p+2+pv)dt+H^f(w,\cdot)+(1+v)A .
\]
At fixed profiles all explicit coefficients and their needed
derivatives are bounded on \(t\ge-\epsilon\).
The derivatives in the axial Hessian direction carry \(d\) or \(\chi\),
as in the proof of Lemma~\ref{lem:ah-floor}; the same holds for
\(p_tdt\), because
\(\partial_p((4+p)/(4+pv))=4d/(4+pv)^2\).
The derivative of \(A\) in the drift \(\chi A(w)w\) is contracted
axially and retains \(\chi\). Coercivity of
Equation~\eqref{eq:ah-T} and Young's Inequality consequently give
\[
 u^{-1}|\operatorname{div}(u\mathcal D)|
 \le\eta\mathcal T+C_N,\qquad
 4|\langle A,dt\rangle_{A_d}|
 \le\eta\mathcal T+C_N
\]
for any fixed \(\eta>0\). These statements bound contracted
expressions; they do not assert a bound for the unweighted axial
entry of \(H^f\).

On a sufficiently high \(Z\) level, either \(t\) is large or \(D\) is large.
In the first case \(3(e^{4t}-1)\) absorbs the bounded errors.
In the second case
\[
 \tau A_Ra\sigma
 =\frac{\tau A_R}{\lambda l_0}\frac{v}{\sqrt d}
\]
does so. At fixed profiles \(\tau A_R/\lambda\) has a positive
lower bound on the whole exterior. The floor penalty is bounded
by a fixed-profile constant and is absorbed in the same manner.
Choosing \(\eta\) below a fixed fraction of \(\delta_0\) yields
\begin{equation}\label{eq:ah-high-divergence}
 \operatorname{div}(4uA_\chi\nabla Z)
 \ge cu(1+\mathcal T+\tau A_Ra\sigma)
 \quad\hbox{on } \{Z>M_N\}
\end{equation}
in the first three stages. The positive \(c\) is fixed; the level
\(M_N\) may depend arbitrarily on all fixed profiles.

\paragraph{The boundary bound before graph comparison.}
In stage one the compatible face equations give
\(-H+w_\nu(F-P_B)=-(1-a)H=O(d)\).
The normal drift has order \(\chi\), since \(A=0\) on the collars.
Thus Equation~\eqref{eq:ah-homotopy-boundary} gives a principal
normal flux of absolute value at most \(Cud\).
The signed slope bound gives
\[
 ud=\lambda L\sqrt d\le C(\tau/\ell)L
\]
on these fixed collars. Multiply
Equation~\eqref{eq:ah-high-divergence} by an increasing cutoff
\(j_0(Z)\), zero up to \(M_N\) and one above \(M_N+1\).
The outer boundary contributes zero. Apply
Equation~\eqref{eq:ah-polynomial-trace} to the inner boundary
term. Its nondifferentiated part is at most
\[
 CN^m\frac{\tau}{\ell}
 \int_{\mathcal C}u j_0(1+\sqrt{\mathcal T}),
\]
which is absorbed by the positive bulk if \(B_0>4m\).
The derivative term is absorbed by
\(\int u j_0'|dZ|_{A_\chi}^2\), leaving a bounded compact strip
integral. We obtain
\begin{equation}\label{eq:ah-high-energy}
 \int u j_0(1+\mathcal T+\tau A_Ra\sigma)
       +\int u j_0'|dZ|_{A_\chi}^2\le C_N .
\end{equation}
Only the polynomial native-weight trace constant was used in
this smallness choice. In particular no comparison involving
\(\ell^{-1}\), no Schauder constant, and no power of
\(n_\infty\) enters the inequality \(B_0>4m\).

\paragraph{The local bound near inner faces.}
After Equation~\eqref{eq:ah-high-energy} is established, constants
may depend arbitrarily on the fixed profiles.
On the ordinary graph \(\Gamma=g+df^2\), the measures
\(dV_\Gamma\) and \(\sqrt D\,dV_g\), and their inverse gradient
metrics, are comparable on a fixed compact patch. Moreover
\[
 e^{2Z}=e^{2t}D^{1/4},\qquad
 D^{1/4}\le C_N(1+\tau A_Ra\sigma).
\]
Thus Equation~\eqref{eq:ah-high-energy}, including the bounded
low strip, gives a local \(L^2(dV_\Gamma)\) bound for \(e^Z\).
The graph mean curvature is bounded by
\(C_N(1+\sqrt{\mathcal T})\): its axial Hessian coefficient is
\((1+|df|^2)^{-1}\le C_Nd\), while Equation~\eqref{eq:ah-T}
controls the transverse and \(d\)-weighted axial entries.
The graph Sobolev argument of Lemma~\ref{lem:fc-graph-iteration}
therefore applies, including its boundary term by
Equation~\eqref{eq:ah-polynomial-trace}. It gives
\[
 \|\omega\|_{L^6(dV_\Gamma)}^2
 \le C_N\int_\Gamma
 \{|\nabla_\Gamma\omega|^2+(1+\mathcal T)\omega^2\}.
\]

For explicit control of the scalar test, multiply its first-stage
equation by \(\eta^2e^{2kZ}/L\), with compact patch cutoff \(\eta\).
The derivative of \(L\) contributes
\((p+2)dt+A-s\). Its first term costs a small \(\mathcal T\)
fraction plus \(C_N|dZ|_{A_\chi}^2\); the spatial term is bounded
at fixed profiles and has the same Young bound. To control the drift,
write \(b_D=K(w,\cdot)+A(w)w^\flat\), so \(\mathcal D=A_\chi b_D\).
On the fixed compact patch, the exact dual-metric identity gives
\[
 |\mathcal D|_{A_\chi^{-1}}^2=|b_D|_{A_\chi}^2
 \le 2(|K|^2+|A|^2)\le C_N .
\]
Thus \(|\mathcal D\cdot dz|\le C_N|dz|_{A_\chi}\).
For \(z=Z\), Young's Inequality bounds the test's \(2k\,dZ\)
term by \(k|dZ|_{A_\chi}^2+C_Nk\). For \(z=t\), use
\(A_\chi\le A_d\) and coercivity to bound the \((p+2)\,dt\)
term by \(\varepsilon_1\mathcal T+C_{N,\varepsilon_1}\),
with \(\varepsilon_1>0\) chosen sufficiently small.
The cutoff terms obey the same dual-metric estimate.
With \(k\ge k_*(N)\),
\begin{equation}\label{eq:ah-moser-energy}
 \int_\Gamma e^{2kZ}\eta^2
       \{\mathcal T+k|\nabla_\Gamma Z|^2\}
 \le C_Nk\int_\Gamma e^{2kZ}(\eta^2+|d\eta|^2).
\end{equation}
At a face the remaining trace terms have only
\(\sqrt{\mathcal T}\) and \(k|\nabla_\Gamma Z|\) growth.
Young's Inequality absorbs them in the left side of
Equation~\eqref{eq:ah-moser-energy}, at cost \(C_Nk\) on the right.
This step requires no new small factor \(\tau/\ell\).

The Sobolev Inequality iterates \(2k\) to \(6k\) on nested patches.
Interpolation with the bounded \(L^2\) norm closes the iteration:
the larger-patch supremum has exponent \(1-1/k_*<1\), so geometric
radius gaps give a uniform smaller-patch supremum. This includes
inner faces. Equation~\eqref{eq:ah-high-divergence} excludes a
higher interior maximum elsewhere, and the outer value is zero.
Hence \(Z\) is uniformly bounded above in stage one.

In stages two and three there is no drift. The high-level source
excludes an interior maximum. At an inner face the upper \(f\)-slope
bound forces \(t>1\) if \(Z\) is sufficiently large. Then \(j=1\)
and the conormal derivative is zero, contradicting the boundary
point principle. In stage four \(f=0\), and the same argument works at
positive scale; at zero scale \(Z=0\) by its energy identity.
Consequently
\[
 -\epsilon\le t\le Z\le C_N,\qquad D\le C_N
\]
throughout the continuation. On compact sets \(|df|\le C_N\);
on the end the same bound holds for \(\lambda|df|\).

\paragraph{Classical estimates.}
On a hyperbolic unit patch centered at \(x_0\), put
\(c=\lambda(x_0)\) and \(f_{\rm rescaled}=cf\). Also divide the
entire scalar divergence equation, its total flux, and its prescribed
conormal datum by the same constant \(c\). In local coordinates its
leading matrix is then
\[
 \widetilde a=4\sqrt{\det g}\,(u/c)A_\chi.
\]
The ratio \(\lambda/c\) is uniformly bounded above and below on the
patch, with bounded logarithmic derivatives. At fixed capped profiles,
\(u/c\) and \(\chi\) have positive upper and lower bounds on the
range just proved. The divided drift and source have the same
normalization, and division by a constant creates no derivative term.
At fixed profiles the trace coefficients are uniformly
elliptic and bounded on the range just proved, and its normalized
source is bounded. The rank-one estimate of
Theorem~\ref{reg-rank-one} gives a positive H\"older exponent for the
rescaled gradient and
\[
 \int_{B_r}|D^2f_{\rm rescaled}|^2\le C_Nr^{1+\gamma_N}
 \quad(\gamma_N>0).
\]
After \(t=t(x,Z,df)\) is substituted in the scalar equation, its
source is bounded by
\(C_N(1+|DZ|^2+|D^2f_{\rm rescaled}|^2)\).
The drift flux and prescribed inner conormal value are bounded.
Lemma~\ref{reg-natural-growth} therefore gives H\"older continuity
of \(Z\), including the disjoint Dirichlet and conormal faces.
The coupled bootstrap of Proposition~\ref{reg-coupled-bootstrap} applies:
the trace coefficients and source are now H\"older; after solving
the trace equation for two derivatives of \(f\), the scalar equation
has H\"older leading coefficients and natural quadratic growth in
\(DZ\). On a face its normal derivative is a smooth function of
\(x,Z,f,df\), because \(df\) is normal. There is no second derivative
of \(f\) or derivative of \(Z\) in that boundary function. These
checks give every local classical bound, uniformly in truncation.

\paragraph{A trace solution for arbitrary inputs.}
The preceding estimates apply only to prospective fixed points.
To define the compact map, fix a finite truncation and a bounded
input \(Z\in C^{1,b}\), without a floor assumption.
At large \(\sigma=|df|\), the defining relation for \(t\) gives
\[
 t=\frac{4Z-\log(\lambda\sigma)}{n+4}+O(1),\qquad
 \chi\asymp\sigma^{-8/(n+4)},\qquad
 \frac{\sqrt D}{l}\asymp\sigma .
\]
Since \(n\ge N\ge4\), \(\chi\ge c/(1+\sigma)\).
The new coefficient \(A\) has growth \(O(\log(2+\sigma))\), but its
normalized trace contribution has size
\[
 C\sigma\chi|A|\le C_N(1+\sigma).
\]
All other terms have the same linear growth on the bounded height
range. The principal matrix is independent of the height, the
normalized source has strict positive height derivative, and all
expressions are smooth and direction free on bounded ranges.
These are exactly the hypotheses of
Proposition~\ref{prop:fc-trace-solve}. It gives the unique
\(C^{3,b}\) trace solution, continuous in the input and continuation
parameter. Its constants may depend on this finite truncation;
this causes no difficulty in defining the map.

Freeze its coefficients, drift, and source, and solve the linear
mixed problem for a new \(Z\), leaving only its principal gradient
unfrozen. The outer Dirichlet face gives coercivity. The output
is bounded in \(C^{2,b}\) on bounded input sets, hence the map is
compact in \(C^{1,b}\). The set defined by
\(\min t(x,Z,df)>-\epsilon\) and a sufficiently large norm cutoff
is open in the parameter--input product.
Lemma~\ref{lem:ah-floor} and the preceding bounds exclude boundary
fixed points. The endpoint map is zero, and zero is admissible.
Lemma~\ref{lem:fc-moving-degree} therefore gives a physical
fixed point. First take \(N\) sufficiently large for the finite
list of polynomial absorptions and compact separations; then take
each outer radius sufficiently large for all fixed-profile barriers.
The a priori bounds are uniform in that outer radius, so a diagonal
compact limit solves the physical system with \(t\ge-\epsilon\).

Finally, on the physical system,
\[
 w(F)=\tau A_Ra\sigma+h\,w(d\log A_R)+w(C).
\]
Young's Inequality and the end height bound give
\begin{equation}\label{eq:ah-capillary-error}
 |h\,w(d\log A_R)|
 \le o(1)\rho+\frac14\tau A_Ra\sigma .
\end{equation}
Indeed it vanishes before \(R\). Beyond \(R\),
\(\tau A_R/(\lambda l_0)\ge c\tau/R\), so its remainder is at
most \(C(R/\tau)|h|^2\); relative to \(\rho\) it is bounded by
\(CR^{4+\delta-2\beta}/\tau=o(1)\).
Equations~\eqref{eq:ah-curvature}, \eqref{eq:ah-offset-budget},
and \(M_d>10\rho\) imply \(R_{\widehat g}\ge-6\), since
\(d_0\le1/2\). The face identity gives
\[
 H_{\widehat g}=e^{-2t}\sqrt d(H+4\partial_\nu t)
             =-e^{-2t}\sqrt d\,N_B<0.
\]
Also \(\widehat g\ge e^{-4\epsilon}g\), so it is complete including
its boundary.
\end{proof}

\subsection{Decay and the four limiting fluxes}
\label{subsec:ah-end}

\begin{lemma}[Differentiated end estimates]\label{lem:ah-end}
For each fixed \(N\), the solution of Proposition~\ref{prop:ah-existence}
satisfies
\[
 \lambda f=O_j(r^{-\beta}),\qquad
 w=O_j(r^{-\beta}),\qquad
 Z-t=O_j(r^{-2\beta})
\]
for every \(j\), and \(Z,t=O_j(r^{-\gamma})\) for every \(\gamma<3\).
Furthermore
\[
 R_{\widehat g}+6=O_j(r^{-3-\delta}).
\]
All four metric fluxes of \(\widehat g\) exist. They satisfy
Equation~\eqref{eq:ah-alg-four-flux}, and the time component satisfies
Equation~\eqref{eq:ah-alg-time-flux}.
\end{lemma}

\begin{proof}
The profiles are fixed throughout this proof. On the ultimate tail
\(n=n_\infty\) is constant and \(A=s\), and \(A_R\) is proportional to
\(r\). The height bound of Lemma~\ref{lem:ah-heights} therefore gives
\(\lambda f=O(r^{-\beta})\). On each hyperbolic unit patch rescale
\(f\) by the central value of \(\lambda\). The trace equation is
uniformly elliptic there by Proposition~\ref{prop:ah-existence}.
Its inhomogeneous term at \(f=0\) is \(O_j(r^{-\kappa})\); the remaining
zeroth-order term is a bounded coefficient times the rescaled \(f\).
The local linear estimates for this fixed smooth solution consequently
give \(\lambda f=O_j(r^{-\beta})\), using \(\beta<\kappa\).
The equations for \(w\) and \(D\) give the next two estimates. The same
argument holds on the outer truncations with uniform boundary estimates
at their zero Dirichlet face.

Here is the scalar expansion before assuming decay of \(Z\).
Since \(t\) is bounded and \(n_\infty\) is fixed, the substitution
\(t=Z+O_j(r^{-2\beta})\) has uniformly bounded differentiated
coefficients. In the flux,
\[
 A_\chi=I+O_j(r^{-2\beta}),\qquad
 u^{-1}du=s+(p+2)dZ+O_j(r^{-2\beta}).
\]
The tensor part of \(S\) has order \(O_j(r^{-\beta})\). In
Equation~\eqref{eq:ah-T}, the part quadratic in \(dZ\) is
\(4(p+1)|dZ|^2\) to error \(O_j(r^{-2\beta})\).
Every mixed term between \(S\) and \(dt\) has an additional factor
\(a\): its two possible forms are \(a\langle M,y\rangle\) and
\(a q_1x_t\). Thus these terms, including their differentiated
bounds, have order \(r^{-2\beta}\). The tensor drift \(K(w,\cdot)\)
has order \(r^{-\kappa-\beta}\).
The leading lapse term \(4\langle s,dZ\rangle\) from
\(u^{-1}\operatorname{div}(4u\nabla Z)\) cancels the explicit lapse
term on the right side of the scalar equation.

The capillary term is \(O_j(r^{-2\beta})\). Even before its support
disappears, the penalty contributes
\(O_j(r^{-3-\delta})+O_j(r^{-2\beta})\), because \(n_\infty\) is
fixed and all derivatives of \(t\) are bounded. Since
\(2\beta>4+\delta\), the resulting equation is
\begin{equation}\label{eq:ah-end-scalar}
 \Delta_gZ+C_N(Z)|dZ|^2
       =\frac34(e^{4Z}-1)+O_j(r^{-3-\delta}),\qquad
 C_N(z)=p(z)+2-d_0(z)(p(z)+1).
\end{equation}
Here \(p(z)=n_\infty p_*(n_\infty z)\) and \(d_0(z)\) denotes the
corresponding composed coefficient. This expansion requires only the
bounded differentiated range already proved.

Define the increasing function
\[
 \Phi(z)=\int_0^z\exp\!\left(\int_0^\xi C_N(\eta)\,d\eta\right)d\xi .
\]
On the bounded solution range both \(\Phi'\) and its reciprocal are
bounded. Writing \(Y=\Phi(Z)\), Equation~\eqref{eq:ah-end-scalar}
becomes
\[
 (\Delta_g-c(Z))Y=O_j(r^{-3-\delta}),\qquad
 c(z)=\frac{3(e^{4z}-1)\Phi'(z)}{4\Phi(z)},\qquad c(0)=3.
\]
The quotient is positive and bounded away from zero on this range.
For \(\gamma=1\),
\(\Delta_g r^{-\gamma}=-r^{-\gamma}+o(r^{-\gamma})\);
hence a large multiple of \(r^{-1}\) bounds both signs of \(Y\)
by the maximum principle, first on truncations whose outer value is
zero and then in their limit. In particular \(Z\to0\), so \(c(Z)\to3\).
For any \(0<\gamma<3\),
\[
 (\Delta_g-c(Z))r^{-\gamma}
   =(\gamma^2-2\gamma-3+o(1))r^{-\gamma},
\]
whose leading coefficient is negative. The same comparison proves
\(Y,Z,t=O(r^{-\gamma})\). Unit-patch estimates for the transformed
equation, followed by differentiation and the already bounded coupled
coefficients, prove the differentiated versions.

Choose \(\gamma>(3+\delta)/2\).
Lemma~\ref{lem:ah-alg-scalar-decay} now applies to
Equation~\eqref{eq:ah-end-scalar} and gives the scalar-curvature
estimate, also after differentiation.
In particular the penalty vanishes sufficiently far out.
Proposition~\ref{prop:ah-alg-mass-passage} applies with
\(s=\gamma>3/2\): all the exponent inequalities there follow from
\(\kappa>2\) and \(\beta>2\). This proves existence of the four limits
and both stated flux identities.
\end{proof}

\subsection{The weighted divergence estimate}
\label{subsec:ah-mass-estimate}

\begin{proposition}[Vanishing loss in the time flux]
\label{prop:ah-mass-loss}
The profiles can be chosen in the stated order so that the physical
solutions satisfy
\[
 \lim_{r\to\infty}\int_{S_r}g(V_*,\nu_g)\,dA_g\ge-o_N(1).
\]
Consequently \(p_0(\widehat g)\le p_0(g)+o_N(1)\).
The error depends on the fixed prepared data and tends to zero as
\(N\to\infty\).
\end{proposition}

\begin{proof}
Use the primitive of Proposition~\ref{prop:ah-alg-primitive}.
For brevity write
\[
 \mathfrak C=\tau A_Ra\sigma
     =\frac{\tau A_R}{\lambda l_0}\frac{v}{\sqrt d},
 \qquad
 b_c=\frac{1-e^{-2t}}2\quad(t\ge0).
\]
The exact divergence is
\begin{equation}\label{eq:ah-mass-divergence}
 u^{-1}\operatorname{div}V_*
   =d_0(\mathcal T+\mathfrak C)+\rho
      -m_0n^{P_0}(\rho_0+v)+\mathcal R,
\end{equation}
where \(\mathcal R\) is Equation~\eqref{eq:ah-alg-primitive-residual}.
We estimate each of its terms. Every Young constant below depends
only on a chosen fixed fraction of the positive terms; it does not
depend on \(N,R\), or a classical bound for the solution.

On \(r\le2R\) first set aside the additive term
\(4\Delta G_c\) in the ordinary-volume divergence. The bounded
logarithmic derivatives of \(n,\lambda\) give
\[
 |\Delta G_c|\le C\lambda/n,\qquad
 \int_{r\le2R}|\Delta G_c|\,dV_g
       \le\frac C N\left(1+\int_1^{2R}r^{2-q}\,dr\right)
       \le\frac C N .
\]
Also, uniformly in this whole region,
\begin{equation}\label{eq:ah-weight-smallness}
 \frac{(1+\log N+\log(1+r))^2}{n\rho}=o_N(1),
\end{equation}
since \(q>3+\delta\). This estimate includes the capped transition,
where \(n\asymp NR^q\).

\paragraph{Positive values beyond the transition in \(t\), \(r\le2R\).}
If \(t\ge1/n\), then \(p=0,A=s\). Integrating the fixed profile on
\(0\le nt\le1\) and differentiating in \(\log n\) through order two
gives
\begin{equation}\label{eq:ah-positive-primitive}
 b_1E_G=b_cE+O(1/n),\qquad
 b_1s_G=b_cs+O(1/n).
\end{equation}
For example the primitive above \(1/n\) is
\(\lambda\{L_\infty(e^{2t}-1)/2+c(n)\}\),
where \(L_\infty>0\) is independent of \(n\) and
\(|(n\partial_n)^jc(n)|\le C_j/n\) for \(j\le2\).
Division by \(\lambda L_\infty e^{2t}\) proves these bounds,
including their uniformity for arbitrarily large \(t\).

Put \(z=e^{2t}\ge1\). Direct factorization gives
\[
 3(e^{4t}-1)-12b_c-36b_c^2
  =\frac{3(z-1)^3(z+3)}{z^2}\ge0.
\]
Since \(E\le3\), the zero-derivative terms supply \(36b_c^2\).
The principal \(F\) error obeys
\[
 4b_c|s(w)F|\le18b_c^2+\frac29F^2.
\]
By Lemma~\ref{lem:ah-alg-coercivity} and \(d_0=1/2\),
the available \(F^2\) coefficient is at least \(1/3\).
Thus strict fractions remain in both places.
The \(K\) term is at most \(Cb_c|s||K|\), and is bounded by a
small fraction of the remaining \(b_c^2\), plus
\(C|s|^2|K|^2\). The lapse smallness in
Lemma~\ref{lem:ah-lapse} makes the latter an arbitrarily small
fraction of \(\rho\).
The errors in Equation~\eqref{eq:ah-positive-primitive} are bounded
by a further small fraction of \(\mathcal T\), plus \(C/n\);
the terms containing \(p\) are zero. Equation~\eqref{eq:ah-weight-smallness}
absorbs this remainder in \(\rho\).

\paragraph{The nonnegative transition band, \(r\le2R\).}
For \(0\le t\le1/n\),
\[
 b_1=O(1/n),\qquad |s_G|+|E_G|\le C,\qquad 0\le p\le n.
\]
The \(F\) term costs \(\eta F^2+C_\eta/n^2\); the \(K\) term
costs \(C/n\) since the prepared \(K\) is bounded.
For the axial and transverse derivative terms, respectively,
\[
 b_1p\,v\sqrt d\,|X_1|
       \le\eta(p+1)dX_1^2+C_\eta b_1^2p,\qquad
 b_1p\,v|y|
       \le\eta(p+1)|y|^2+C_\eta b_1^2p .
\]
Their remainders are \(O(1/n)\).
The zero-derivative negative term has the same order; the
cosmological term is nonnegative. Coercivity and
Equation~\eqref{eq:ah-weight-smallness} absorb every cost.

\paragraph{Negative \(t\), \(r\le2R\).}
Write \(X=-nt\ge0\). Then
\[
 G=\frac{\lambda e^{(n+2)t}}{n+2},\qquad
 b_1=\frac1{n+2},\qquad
 |s_G|\le C(1+X),\qquad |E_G|\le C(1+X)^2 .
\]
The first two derivative bounds follow by differentiating
\(\log G=\log\lambda+(n+2)t-\log(n+2)\) with \(t\) fixed.
The same Young estimates as above, now with \(p=n\), show that all
terms other than a small fixed fraction of \(\mathcal T\) cost
\[
 C_\eta\frac{(1+X)^2}{n}.
\]
This includes the cosmological term, since
\(1-e^{4t}\le4|t|=4X/n\).
For \(X\le C_0(\log N+\log(1+r))\), Equation~\eqref{eq:ah-weight-smallness}
absorbs the cost in \(\rho\).

Choose the fixed \(C_0\) larger than the exponents in
\(1/\tau,\lambda\), and the polynomial \(X^2\). For larger \(X\),
if \(v\ge1/2\), the ratio of this cost to \(\mathfrak C\) is bounded
by
\[
 C\frac{\lambda(1+X)^2e^{-X}}{n\tau A_R},
\]
which is \(o_N(1)\), uniformly in this region.
If \(v<1/2\), then \(d>1/2\) and
\(u\le C\lambda e^{-X}\). After multiplication by \(u\), the
unabsorbed cost is at most \(C\lambda/n\), and its total integral
on \(r\le2R\) is \(O(1/N)\).
These two alternatives account for all negative \(t\) in this region.

\paragraph{The capped tail, \(r\ge2R,\ t\ge0\).}
Here \(n=n_\infty\), \(s_G=A=s\), and \(E_G=E\le E_0\), where
\(E_0=\Delta\lambda/\lambda\le3\). Set
\[
 \gamma_c=\frac{G_c}{\lambda L}=\frac1{(n+2)L},\qquad
 c=b_1-\gamma_c=\frac{\int_0^tL(z)\,dz}{L(t)}.
\]
Since \(p\ge0\), \(L'/L=p+2\ge2\); therefore \(0\le c\le b_c\).
Also \(L\ge1\), so \(0\le\gamma_c\le1/(n+2)\).
Combining the primitive and correction terms gives
\[
 -4b_1E+4\sqrt d\,\gamma_cE_0
 =-4cE+4\gamma_c(\sqrt d\,E_0-E)
 \ge-12b_c-Cv/n.
\]
For the last inequality use the exact expression for \(E\) to
bound \(|E-E_0|\le Cv\), and \(1-\sqrt d\le v\).

Split the \(F\) coefficient into \(c+\gamma_c\).
The \(c\) term is bounded by
\(18b_c^2+(2/9)F^2\), and the \(\gamma_c\) term costs a small
remaining fraction of \(F^2\), plus \(Cv/n^2\).
The \(K\) term containing \(c\) is absorbed using the lapse smallness
as before. Its \(\gamma_c\) part satisfies
\[
 \gamma_c a|s||K|\le\eta\rho+C_\eta v/n^2.
\]
When \(p=0\) there are no further terms. When \(p>0\), necessarily
\(t<1/n\) and \(b_1=O(1/n)\); the two derivative terms cost
\(\eta\mathcal T+C_\eta v/n\) by the preceding shifted-square estimates.
Finally, on this tail,
\[
 \frac{\tau A_R}{\lambda}\ge c\tau/R,\qquad
 l_0\le L_\infty,\qquad
 \frac1{n_\infty}=o(\tau/R).
\]
Thus every remaining \(Cv/n\) is absorbed by an arbitrarily small
fixed fraction of \(\mathfrak C\), using
Equation~\eqref{eq:ah-small-ratios}.

\paragraph{The capped tail, \(r\ge2R,\ t<0\).}
Put \(x=(n+2)|t|\) and \(X=n|t|\). The zero-gradient potential,
including the centering correction, is
\begin{equation}\label{eq:ah-negative-tail-potential}
 P(t)=3(e^{4t}-1)+\frac{4E_0}{n+2}(e^x-1).
\end{equation}
If \(x\le1\), the inequalities \(e^{-y}\ge1-y\) and
\(e^x\ge1+x+x^2/2\) give
\[
 P(t)\ge\frac{-4(3-E_0)x+2E_0x^2}{n+2}
       \ge-\frac{2(3-E_0)^2}{(n+2)E_0}.
\]
Since \(3-E_0=O(r^{-2})\), this is \(O(r^{-4}/n)\), hence
\(o_N(1)\rho\) because \(\delta<1\).
For \(x\ge1\),
\((e^x-1)/x\ge e-1\), so the same linear bound for \(e^{-4|t|}\)
shows \(P(t)\ge0\) as soon as \(E_0(e-1)\ge3\).
This holds on the whole tail for all sufficiently large \(N\).

The difference from the zero-gradient potential is at most
\(Ce^xv/n\), because
\(|E-E_0|+|1-\sqrt d|\le Cv\).
The floor gives \(e^x\le e^{2\epsilon}e^X\le Ce^X\).
The \(F\), axial, and transverse terms, after a small fraction of
\(\mathcal T\), cost \(Cv/n\); the \(K\) term costs a small
fraction of \(\rho\), plus \(Cv/n^2\).
All these residuals are at most \(Ce^Xv/n\), which is absorbed by
\(\mathfrak C=(\tau A_R/\lambda)e^Xv/\sqrt d\).
Again the required ratio is exactly \(R/(n_\infty\tau)=o(1)\).

We have proved that, away from the penalty, the integral of
Equation~\eqref{eq:ah-mass-divergence} retains positive fixed
fractions of \(u\rho,u\mathcal T,u\mathfrak C\), up to an ordinary
volume error with total integral \(o_N(1)\).

\paragraph{The penalty band.}
On its support,
\[
 -\epsilon\le t\le-\epsilon+1/n,\qquad
 l_0\asymp e^{-\epsilon n},\qquad e^{2t}\asymp1,
\]
with constants independent of \(N\).
On \(n<2N\), which is a fixed compact region, \(\lambda,A_R\)
and the weights are bounded, and \(l_0\le C\ell\).
With a fixed polynomial \(P(N)\), the definitions imply
\[
 u\,m_0n^{P_0}(\rho_0+v)
 \le P(N)(\ell+\ell^2\sigma)
 \le\eta\,u\mathfrak C
              +P(N)(\ell+\ell^2/\tau).
\]
For the last inequality use the exact identity
\(u\mathfrak C=e^{2t}\tau A_R l^2\sigma^2\)
and apply Young before replacing \(l_0\) by its upper bound.
Thus the integrated remaining cost on this fixed compact region
is at most
\[
 P(N)(N^{-B_0}+N^{-3B_0/4})=o_N(1)
\]
once \(B_0\) exceeds the fixed polynomial exponents.

On \(n\ge2N\), compare the penalty's \(v\) term directly with
the capillary term. Their quotient is at most
\[
 \frac{\lambda n^{P_0}e^{-\epsilon n}}{\tau A_R}.
\]
This tends to zero uniformly. To see the order without using
\(n_\infty\) as an uncontrolled polynomial constant, put \(y=n/N\ge2\).
Before capping, \(\lambda/A_R\le C y^{1/q}\); after capping the same
bound holds with \(y=n_\infty/N\), since \(\lambda/A_R\asymp R\).
The displayed quotient is therefore bounded by
\[
 C N^{P_0+5B_0/4}\,y^{P_0+1/q}N^{-B_0y}.
\]
For large \(N\) this is decreasing in \(y\ge2\), with maximum
\(O(N^{P_0-3B_0/4})\). Increasing \(B_0\) makes it vanish.
The \(\rho_0\) term is treated the same way where \(v\ge1/2\),
since \(\rho_0\) is bounded. Where \(v<1/2\), use \(d>1/2\):
its ordinary-volume cost is at most
\[
 C\lambda\rho_0\, n^{P_0}e^{-\epsilon n}.
\]
The supremum of its final factor for \(n\ge2N\) tends to zero,
and \(\int\lambda\rho_0\,dV_g<\infty\).
This proves the required \(o_N(1)\) total penalty loss.

\paragraph{The inner boundary.}
The correction leaves \(V_*=V\) on every inner collar.
On a physical face,
\[
 V_\nu/u=-(1-a)H-N_Bd\ge-Cd,\qquad
 \sqrt d\le C\tau/\ell .
\]
Thus the negative boundary integral is at most
\(C(\tau/\ell)\int_{\partial\Omega}L\).
Equation~\eqref{eq:ah-polynomial-trace}, with a constant test and
a fixed collar cutoff, bounds it by
\[
 P(N)\frac{\tau}{\ell}
       \int_{\mathcal C}u(\rho+\delta_0\mathcal T)\,dV_g.
\]
Here \(\rho\) has a fixed positive minimum on the collars.
Since \(\tau/\ell=N^{-B_0/4}\), choosing \(B_0\) larger than these
fixed polynomial exponents absorbs this loss in the positive
fractions already retained.

The parameter order is now complete: \(\delta_0\) and \(P_0\)
were fixed by the floor; choose \(B_0\) larger than the finitely
many compact polynomial requirements, and then choose \(a\) by
Equation~\eqref{eq:ah-R-conditions}. No estimate used a polynomial
bound for a fixed-profile Schauder constant or for \(n_\infty\).
Apply the divergence theorem on finite exterior domains. With
\(\nu\) pointing into the exterior at inner faces it gives
\[
 \int_{S_r}g(V_*,\nu_g)\,dA_g
 =\int_{\Omega\cap\{r'\le r\}}\operatorname{div}V_*\,dV_g
      +\int_{\partial\Omega}V_\nu\,dA_g .
\]
The preceding estimates give a lower bound \(-o_N(1)\) after
dropping the remaining nonnegative integrals. Lemma~\ref{lem:ah-end}
supplies the limiting outer flux, so the asserted estimate follows.
Its mass consequence has the stated sign by
Equation~\eqref{eq:ah-alg-time-flux}.
\end{proof}

\subsection{The full-cut inequality and the invariant mass}
\label{subsec:ah-conclusion}

\begin{proof}[Proof of Theorem~\ref{thm:ah-component}]
Fix one prepared data set and the compact black--white exterior used
in the construction. Every full enclosing cut in this smaller
exterior is a full enclosing cut in the prepared original exterior.
Since \(t\ge-\epsilon\) and the graph term is nonnegative,
\(\widehat g\ge e^{-4\epsilon}g\). On each tangent two-plane its
area density is therefore at least \(e^{-4\epsilon}\) times the
\(g\)-density. Taking the infimum over all full cuts gives
\[
 A_{\min}(\partial\widehat\Omega;\widehat g)
       \ge e^{-4\epsilon}A_{\min}(S;g).
\]
This comparison uses the entire intrinsic boundary of each cut,
including the original boundary when the exterior itself is chosen.

By Proposition~\ref{prop:ah-existence} and Lemma~\ref{lem:ah-end},
\((\widehat\Omega,\widehat g,0)\) has one complete AH end with the
required tensor-norm decay and weighted scalar-curvature
integrability, all four finite metric fluxes, scalar curvature
at least \(-6\), and strictly mean-negative boundary. Thus the
time-symmetric case of Theorem~\ref{thm:str-one-sign} applies
in the designated chart, irrespective of the causal character
of its flux covector. Together with Proposition~\ref{prop:ah-mass-loss},
it gives
\[
 p_0(g)+o_N(1)\ge p_0(\widehat g)
 \ge
 \sqrt{\frac{e^{-4\epsilon}A_{\min}(S;g)}{16\pi}}
 \left(1+\frac{e^{-4\epsilon}A_{\min}(S;g)}{4\pi}\right).
\]
Send \(N\to\infty\); then \(\epsilon=B_0\log N/N\to0\).
Finally remove the strict preparation using
Lemma~\ref{lem:ah-preparation}, which preserves the four metric
flux limits and the full-cut infimum. This proves
Equation~\eqref{eq:ah-component-bound}.
\end{proof}

\begin{theorem}[The AH mass inequality]\label{thm:ah-invariant}
If the metric flux covector of the data in
Theorem~\ref{thm:ah-component} is future timelike, then
\[
 m_{\rm AH}:=\sqrt{p_0^2-p_1^2-p_2^2-p_3^2}
 \ge
 \sqrt{\frac{A_{\min}(S;g)}{16\pi}}
 \left(1+\frac{A_{\min}(S;g)}{4\pi}\right).
\]
\end{theorem}

\begin{proof}
The four static potentials form the defining Lorentz representation
of the orientation and time-orientation preserving hyperbolic
isometry group. The metric flux is linear in the potential.
Consequently, composing the designated end chart with such an
isometry acts on its four components by the corresponding Lorentz
transformation. This follows directly by changing variables in
the flux formula: the background metric and its connection are
preserved, and the replacement sphere exhaustion has the same
limit because the weighted scalar-linearization divergence is
absolutely integrable. The quadratic remainder is integrable
under the stated decay exponent \(>3/2\).
The tensor decay and weighted integrability conditions persist,
since two radial functions related by a fixed hyperbolic isometry
are comparable far out.

A future timelike covector admits a proper future Lorentz
transformation taking it to \((m_{\rm AH},0,0,0)\).
Apply Theorem~\ref{thm:ah-component} in this chart. The geometry of
the exterior, its future boundary signs, and its intrinsic full-cut
area infimum have not changed. This proves the assertion.
Only a hyperbolic change of spatial end coordinates was used.
\end{proof}

This completes the proof of Theorem~\ref{thm:main-numerical}.
